% Автоматически перенесено из ИСТОРИОГРАФИЯ_ВЕРСИЯ_3.docx

\section*{Историография}
\addcontentsline{toc}{section}{Историография}

\subsection*{§1. Западноевропейская медиевистика о каталонском обычном праве XII--XIII~вв.}
\addcontentsline{toc}{subsection}{§1. Западноевропейская медиевистика о каталонском обычном праве XII--XIII~вв.}

Параграф посвящён становлению и развитию западноевропейской традиции изучения каталонского обычного права. Особый интерес представляет смена исследовательских подходов --- от позитивистского описания источников к проблемно-ориентированным и междисциплинарным исследованиям обычного права. В этом параграфе мы рассмотрим, в каком историографическом контексте сформировались современные представления каталонском обычном праве XII--XIII~вв.

\subsubsection*{§1.1. Ранний этап изучения каталонского обычного права XII--XIII~вв.}
\addcontentsline{toc}{subsubsection}{§1.1. Ранний этап изучения каталонского обычного права XII--XIII~вв.}

В центре внимания находятся ранние исследования, посвящённые правовой культуре Каталонии XII--XIII~вв., сформировавшиеся в русле правовой истории и романской филологии. Проанализируем их источниковую базу, исследовательские цели и характер интерпретации обычного права. Особое внимание уделим тому, как ранние авторы соотносили каталонский материал с общими схемами развития средневекового права в Европе.

Анализ целесообразно начать с работ историко-правового направления конца XIX~в. --- 1930-х гг., опуская детальное обсуждение первых сочинений, тщательно проанализированных Г.М.~Броком-и-де-Монтагутом\footnote{\emph{Brocà i de Montagut G.M. de.} Els Usatges de Barcelona // Anuario del Instituto de Estudios Catalanes. Barcelona: Institut d'Estudis Catalans, 1915. P.~357--389.}. Для раннего этапа характерен интерес исследователей к «Барселонским обычаям», интерпретируемым многими как «основной закон» средневековой Каталонии без учёта более широкой проблематики.

Большинство исследователей конца XIX -- начала XX~вв. датировали составление этого правового памятника временем правления Рамона Беренгера ~I Старого (1036--1076~гг.), опираясь либо непосредственно на позднесредневековую традицию, либо на аргументы, утратившие убедительность в современном исследовательском контексте (одним из примеров является ссылка Дж.~Королеу-и-Инглады на мавристский труд «Искусство проверять даты исторических событий»\footnote{\emph{Coroleu i Inglada~J. }Código de los Usajes de Barcelona. Estudio crítico // Boletín de la Real Academia de la Historia. Vol.~4. Cuad. II. Madrid, 1884. P.~85--104.}). На её сомнительность уже указывал В.К.~Пискорский\footnote{\emph{Пискорский}\emph{~}\emph{В.К}. Крепостное право в Каталонии в Средние века. Киев: Тип. Ун-та св. \emph{Владимира Н.Т.}~Корчак-Новицкого, 1901. C.~5.}, подходивший к тексту «Барселонских обычаев» как историк, однако, переосмысление их значения в средневекового каталонского права тогда ещё не становилось предметом специального анализа.

Ф.~Фита-и-Коломэ следует позитивистской исследовательской традиции, сосредоточенной на выявлении, введении в научный оборот и текстологическом анализе правовых источников, например, «Барселонских обычаев»\footnote{\emph{Fita y Colomé F.} Cortes y usajes de Barcelona en 1064. Textos inéditos // Boletín de la Real Academia de la Historia. 1890. T.~5. №~17. P.~385--428; \emph{Idem.} El obispo Guisliberto y los Usajes de Barcelona // Boletín de la Real Academia de la Historia. 1891. T.~6. №~18. P.~228--246.}. Основное внимание исследователь уделял установлению подлинности, датировки и авторства текстов, реконструкции их происхождения и связи с политической историей.

Среди работ, посвящённых другим локальным записям обычного права в Каталонии указанного периода, стоит назвать четырёхтомную публикацию текста «Обычаев Тортосы» Б.~Оливера-и-Эстрейера, снабжённую вступительной статье и обширным историческим комментарием\footnote{Costums de Tortosa; Libre de les costums generals scrites de la insigne ciutat de Tortosa: texto auténtico de siglo XIII que da a luz nuevamente, ilustrado con observaciones crítico-literarias y un copioso vocabulario / ed. B.~Oliver y Esteller. Madrid: Impr. de M. Ginesta, 1881.}.

А.-Ш.~Жиро, О.~Брютайльс и Ю.~Фикер в своих изысканиях опирались преимущественно на поздние латинские рукописи и печатные издания «Барселонских обычаев» и иных памятников интересующего нас периода. С одной стороны, это позволяло им рассматривать обычное право в широком историко-правовом контексте. С другой --- приводило к сужению исследовательской перспективы до уровня текстологической критики и поиска «источников» конкретных норм.

Показательным примером такого подхода стала работа Ю.~Фикера, в которой он выявлял заимствования в составе «Барселонских обычаев и отмечал слабые места в первом издании текста А.-Ш.~Жиро\footnote{\emph{Giraud}\emph{~}\emph{Ch.} Op. cit. P.~465--509.}, однако понимал каталонское обычное право исключительно как результат рецепции римско-канонической доктрины, а не как продукт специфической правовой эволюции графства Барселонского\footnote{\emph{Ficker~J.} Über die Usatici Barchinonae und deren Zusammenhang mit den Exceptiones Legum Romanorum // Mitteilungen des Instituts für Österreichische Geschichtsforschung. 1888. Erg.-Bd.~II. S. 236--276.}. Аналогично О.~Брютайльс в статье, посвящённой анализу 72-й статьи «Барселонских обычаев», предложил детальный юридический комментарий, но трактовал норму как отвлечённое правовое понятие, практически проигнорировав исторический контекст, с которым было связано составление данного памятника\footnote{\emph{Brutails}\emph{~}\emph{A.} Étude sur l’article 72 des Usages de Barcelone: Connu, sous le nom cle loi stratæ // Nouvelle revue historique de droit français et étranger. 1888. Vol.~12. P.~59--79.}.

В первой трети XX~в. историография продемонстрировала усиление интереса к каталонскому обычному праву в целом. Итальянский историк права Е.~Беста, задавшись опросом о статусе «Барселонских обычаев» в системе права графства Барселонского, поставил проблему иерархии норм и соотношения писанного и неписаного права, графской правотворческой практики и местных обычаев\footnote{\emph{Besta~E.} Usatici ed usi curiali di Barcellona // Rendication Istituto Lombardo. 1925. Ser. II, Vol.~58. P.~637--652.}.

Исследователи политико-правовых институтов Каталонии, такие как Ф.~Вальс-и-Табернер\footnote{\emph{Valls i Taberner F.} El problema de la formació dels «Usatges» de Barcelona // Revista de Catalunya. 1925. №~7. P.~26--33; \emph{Idem.} Les consuetuds i franqueses de Barcelona de 1284, o «Recognoverunt proceres» // Revista de Catalunya. 1927. №~39. P.~248--254; \emph{Idem.} Els elements fonamentals del Dret català antic: Conferència pronunciada a l’Ateneu Barcelonès, el dia 7 de febrer de 1928 // Revista de Catalunya. 1928. №~47. P.~467--479; \emph{Idem.} La cour comtale Barcelonaise // Revue historique de droit français et étranger. 1935. №~4 (14). P.~662--682; \emph{Idem.} El problema de la formació deIs Usatges de Barcelona // Revista de Catalunya. 1925. №~7 (2). P.~26--33; \emph{Idem.} Els «Usualia de curialibus usibus Barchinonae» (Assaig de reconstrucció) // Estudis universitaris catalans. 1934. №~19. P.~270--280; \emph{Idem.} Noves recerques sobre els Usatges de Barcelona // Estudis universitaris catalans. 1935. №~20. P.~69--83.}, Р.~д’Абадаль-и-Виньялис\footnote{\emph{d’Abadal Vinyals R. de.} La Data i el lloc de la mort del comte Berenguer Ramon I // Butlletí de la Societat Catalana d’Estudis Històrics. 1952. №~1. P.~43--44; \emph{Idem.} El renaixement monàstic a Catalunya després de l’expulsió dels sarraïns // Studia monastica. 1961. №~3. P.~165--177. \emph{Idem.} La institució comtal carolíngia en la pre-Catalunya del segle IX // Anuario de estudios medievales. 1964. №~1. P.~29--75; \emph{Idem.} El dominio Carolingio en la marca peninsular hispánica, siglos IX y~X // Hispania: Revista española de historia. 1968. №~2. P.~33--50.}, Дж.Дж.~Перманьер-и-Айятс, а также итальянский историк права К.Г.~Мор\footnote{\emph{Mor C.G.} En torno a la formación del texto de los «Usatici Barchinonae» // Anuario de Historia del Derecho Español. 1958. T.~27--28. P.~413--460.}, также расширили тематическое поле исследований. Изучая генезис графской и королевской власти, кортесов, муниципального самоуправления и феодальных связей, они неизбежно касались того, как обычай, договор и судебная практика структурировали социальные отношения. Так, Ф.~Вальс-и-Табернер и К.Г.~Мор, анализируя эволюцию институтов Арагонской Короны и каталонских кортесов, показали роль локальных обычаев в формировании городских корпораций. Р.~д’Абадаль-и-Виньялис и Б.~Оливер-и-Эстрейер исследовали особенности взаимодействия власти графов-королей, местной знати и должностных лиц\footnote{\emph{Oliver y Esteller B.} Historia del derecho en Cataluña, Mallorca y Valencia. Código de las costumbres de Tortosa. Madrid: Impr. de M. Ginesta, 1876. 4 vol.}.

Хотя в работах этих авторов изучение «Барселонских обычаев» ещё во многом следовало сложившейся исследовательской традиции, в том числе в вопросе датировки памятника XI~в., как, например, у М.А.~Марина-и-Луны\footnote{\emph{Marín i Luna M.A.} Els Usatges de Barcelona // Revista Jurídica de Catalunya. Barcelona, 1930. Vol.~XXXVI (Abril-maig-juny 1930). P.~97--135.}. Именно последователи классического историко-правового подхода впервые обратили внимание на социальные практики, стоящие за нормами, и на то, как язык закрепляет отношения и компромиссы внутри каталонского общества XII--XIII~вв.

Одним из первых таких исследований стала работа испанского исследователя Д.Дж.~Балари-и-Джовани «Начала каталонской истории»\footnote{\emph{Balari y Jovany D.J.} Origenes Históricos de Cataluña. Barcelona: Hijos de J. Jepús, 1899.}, и в особенности, её параграф, посвящённый «Барселонским обычаям»\footnote{\emph{Idem.}~Op. cit. P.~453--466.}. Это исследование было одной из первых попыток междисциплинарного анализа «Барселонских обычаев», их структуры, содержания и роли в истории Каталонии. Разделяя статьи обычаев на группы, он пришёл к выводу о постепенном складывании текста «Барселонских обычаев» к середине XII~в., то есть времени правления Рамона Беренгера IV или его сына Альфонсо II.

Однако в отличие от других исследователей своего времени Д.Дж.~Балари-и-Джовани не прибегал к толкованию и анализу отдельных правовых установлений, вошедших в текст «Барселонских обычаев». Для него определение датировки составления текста обычаев было лишь способом реконструировать политическую историю графства.

Подводя итог, отметим, что ранние исследования были преимущественно описательными и компилятивными, а их авторы фиксировали существование каталонского обычного права и его особенности, но редко ставили перед собой аналитические задачи и выходили за пределы текстологических исследований.

Зачастую каталонский материал привлекался в качестве иллюстрации к более широким моделям эволюции феодального права, что приводило к сглаживанию его специфики. В то же время именно в этих работах был введен в научный оборот ряд ключевых рукописей и публикаций «Барселонских обычаев», «Обычаев Тортосы», а также Жироны, Льейды, Миравета и др. городских сообществ. Несмотря на методологическую ограниченность, ранний этап заложил источниковедческую основу, без которой последующие исследования были бы невозможны.

\subsubsection*{§1.2. Классические исторические исследования по истории обычного права в Каталонии}
\addcontentsline{toc}{subsubsection}{§1.2. Классические исторические исследования по истории обычного права в Каталонии}

Во 2-й пол. XX~в. сформировалось направление историко-правовых исследований, связанных с развитием теории «феодальной революции». Оно располагалось на пересечении социальной, политической и правовой истории, что подтверждается содержательными историографическими обзорами последних десятилетий\footnote{\emph{Sabaté}\emph{ }\emph{i}\emph{ }\emph{Curull}\emph{ F.} The Catalonia of the 10\textsuperscript{th} to12\textsuperscript{th }centuries and the historiographic definition of feudalism // Catalan Historical Review. 2010. Vol.~3. P.~31--53.}. Исследования П.~Боннасси, Т.Н.~Биссона и П.Х.~Фридмана задали объяснительную модель, связывающую формирование городских обычаев с практиками перераспределения юрисдикций и формированием судебных инстанций.

Обзоры Ф.~Сабате-и-Куруля уточнили понятие феодализма применительно к каталонскому материалу, выделив региональные различия, длительность процессов и место обычая во взаимоотношениях власти и общества. Панорама, предложенная Дж.М.~Сальраком-и-Маресом для каталонской историографии 1970--1990~гг., зафиксировала поворот каталонской медиевистики к структурам власти, письменности и институциональной истории, что создало методологическую базу для изучения локального обычая и судебной практики\footnote{\emph{Salrach J.M.} La historiografía medieval en Cataluña, 1970--1990 // Cuadernos de Historia de España. 2022. №~89. P.~113--154. DOI: 10.34096/che.n89.12228. URL: \url{https://revistascientificas.filo.uba.ar/index.php/che/article/view/12228} (дата обращения: 01.05.2026).}. Эти обзоры, вкупе с общими работами по социально-экономической и политической истории Каталонии, сместили фокус от «писаного» публичного права к повседневным правовым практикам, в рамках которых в том числе и бытовало обычное право XII--XIII~вв.

Важное место в историографии правовой культуры занимает творчество Дж.М.~Понса Гури, который, опираясь на богатый круг локальных источников, реконструировал формы бытования обычного права в Жироне и её окрýге. В работах о жиронских сборниках права он показал, как в условиях сосуществования эмфитевзиса\footnote{\textbf{Эмфитевзис} (лат. emphyteusis от греч. {\unicodefallback ἐ}μφύτευσις) --- постоянная наследственная аренда полевых земельных участков. (см. \emph{Бартошек}\emph{ М.} Римское право: понятия, термины, определения / пер. с чеш. \emph{Ю.В.~Преснякова}; спец. науч. ред. и авт. предисл. \emph{З.М.~Черниловский}. М.: Юридическая литература, 1989. С.~120).} и феодальных земельных держаний местные обычаи выступали инструментом упорядочения отношений крестьян и их сеньоров\footnote{\emph{Pons}\emph{ }\emph{Guri}\emph{ }\emph{J}\emph{.}\emph{M}\emph{.} Entre l’emfiteusi i el feudalisme (els reculls de dret gironins) // Estudi general: Revista de la Facultat de Lletres de la Universitat de Girona. Nº 5--6. 1985--1986. P.~411--418.}. Изучив «дурные обычаи» на территории вокруг Жироны, Дж.М.~Понс~Гури обратился к актам и судебным материалам\footnote{\emph{Idem.} Relació jurídica de la remença i els mals usos a les terres gironines // Revista de Girona. 1986. Vol.~32. P.~36--43.}. Особое значение для настоящего исследования имеют его источниковедческие изыскания, посвящённые раннесредневековому праву и документам об использовании «Барселонских обычаев» до XIII~в. и позволяющие проследить, как общекаталонские нормы интерпретировались в локальной среде\footnote{\emph{Idem.} Documents sobre aplicació dels Usatges de Barcelona anteriors al segle XIII // Acta historica et archaeologica mediaevalia. 1993. №~14--15. P.~39--46.}.

Вместе с тем в работах о каталонских сборниках обычного права XIII--XVIII~вв. он выделил их региональные особенности, что позволило увидеть длительную траекторию формирования и кодификации местных обычаев XII--XIII~вв.\footnote{\emph{Idem., Pladevall Font A.} Particularismes catalans en els costumaris deis segles XIII--XVIII // II Congrés litúrgic de Montserrat. 1967. №~3 (2). P.~103--159.}. Для изучения правовой культуры Жироны принципиально важно, что Дж.М.~Понс~Гури рассмотрел эти источники не только как нормативные тексты, но и как свидетельство правосознания и повседневных стратегий различных социальных групп\footnote{\emph{Idem. }Document d’aplicació del dret senyorial // Recull d'estudis d'història jurídica catalana. 1989. T.~2. P.~193--196.}. Его подход, сочетавший текстологический анализ сборников обычаев с реконструкцией судебной практики, задал важный методологический ориентир для источниковедческого исследования локального обычного права, предпринятого в настоящей диссертации.

Один из наиболее авторитетных исследователей каталонской истории среди медиевистов П.~Боннасси работал в русле французской школы «Анналов», соединив социально-экономическую историю с анализом массовых источников (по большей части архивных и неопубликованных) и правовых обычаев. Фундаментальное значение имели его исследования, прежде всего двухтомная монография «Каталония с середины X~в. до конца XI~в.: рост и трансформация общества»\footnote{\emph{Bonnassie P.} La Catalogne du milieu du Xe à la fin du XIe siècle: croissance et mutations d’une société. Toulouse: Association des publications de l’Université de Toulouse-Le Mirail, 1975. T.~2. P.~718--728.}, в которой правовой блок во многом построен на материале «Барселонских обычаев» в сопоставлении с широкой актовой базой\footnote{\emph{Idem.}~Les conventions féodales dans la Catalogne du XIe siècle // Annales du midi. 1968. Т.~80. P.~529--561.}. Синтез 1970-х гг. закрепил эти наблюдения и привёл исследователя к реконструкции политической структуры принципата, опирающейся на систему верности и сеньориально-вассальные механизмы, выявляемые через «Барселонские обычаи» и актовый материал. В статьях и очерках о каталонском феодализме он обобщил результаты, сохранив акцент на правовых практиках в качестве ядра социальной трансформации\footnote{\emph{Idem.}~El feudalisme català, segle XI // L’avenç. Història dels Països Catalans. 1988. T.~117. P.~24--28.}.

Интерес к коллективным формам самоуправления нашёл отражение в совместной с П.~Гишаром статье о сельских общинах Каталонии и Валенсии, где авторам удалось проследить эволюцию локальных обычаев с IX до середины XIV~вв.\footnote{\emph{Idem., Guichard P.} Les communautés rurales en Catalogne et dans le pays valencien (IXe milieu XIVe siècle). 1984. P.~79--115.} Обзор «Возникновение Каталонии, VIII --- середина XII~вв.» и сравнительный очерк о Тулузском и Барселонском графствах в 801--1213~гг. перенаправили анализ П.~Боннасси в русло исследования “longue durée” и непосредственно вышли за пределы интересующего нас периода XII--XIII~вв.\footnote{\emph{Idem.}~Emergence de la Catalogne, VIIIe --- milieu XIIe siècle // Histoire des Espagnols 1. 1985. P.~159--183; \emph{Idem.}~Le comté de Toulouse et le comté de Barcelone au début du IXe au début du XIIIe siècle: P.~801--1213 -- esquisse d’histoire comparée // Actes del vuitè Col.loqui Internacional de Llengua i Literatura Catalanes. 1989. Т.~1. P.~27--45.} В статье о легендарном происхождении Каталонии автор зафиксировал направления исследования институтов и обычаев--- от публично-властных норм к сеньориальным\footnote{\emph{Idem.}~Sur les «Origines» de la Catalogne: Quelques remarches et orientations de recherche // Symposium internacional sobre els orígens de Catalunya (segles VIII--XI). 1991. Т.~1. P.~437--445.}.

В 1994~г. П.~Боннасси сформулировал тезис о специфике каталонского «процесса феодализации» в сопоставлении с югом Франции, связав судопроизводство и формы ренты с изменениями социальной стратификации\footnote{\emph{Idem.}~El proceso der feudalización en Cataluña y Francia del Sur: similitudes y diferencias // Los orígenes del feudalismo en el mundo mediterráneo. 1994. P.~101--117.}. В одной из поздних статей о «языковых новациях и социально-экономических изменениях» в Каталонии он ввёл лингвистические индикаторы для датировки и интерпретации правовых терминов\footnote{\emph{Idem.}~Nouveautes linguistiques et mutations economico- sociales dans la Catalogne des IXe-XIe siècles // Langages et peuples d’Europe. 2002. P.~47--66.} В целом его синтетические построения, включая реконструкцию политической структуры принципата на основе «Барселонских обычаев», остаются ориентиром для современных исследователей, хотя и требуют уточнения в частных вопросах, например точности датировки составления «Барселонских обычаев»\footnote{Подробнее о биографии П.~Боннасси, логике его творчества и тех, на кого он оказал влияние, см. мемориальный очерк: \emph{Филиппов И.С}. Пьер Боннасси (24.11.1932 --- 14.03.2005) // СВ. 2006. №~67. С.~316--325.}.

Л.То~Фигерас продолжил исследования по истории крестьянства в Каталонии, связав формы зависимости и семейные стратегии с институтами наследования и землепользования. Он уточнил, как из повседневных потребностей возникли относительно устойчивые нормы обычного права. В работах 1990-х~гг. о практиках наследования он реконструировал становление неделимого наследования в X--XII~вв. как механизма концентрации имущества и стабилизации сеньориальной власти\footnote{\emph{To Figueras L.} Señorío y familia: los orígenes del «hereu» catalán (siglos X--XII) // Studia historica. Historia medieval. 1993. №~11. P.~57--80; \emph{Idem.}~Anthroponymie et pratiques successorales (à propos de la Catalogne, Xe--XIIe siècle) // Actes du colloque international organisé par l’École française de Rome avec le concours du GDR 955 du C.N.R.S. «Genèse médiévale de l’anthroponymie moderne» (Rome, 6--8 octobre 1994). 1996. Vol.~226. P.~42--433; \emph{Idem.}~Família i hereu a la Catalunya nord-oriental (segles X--XII). Barcelona: Publicacions de l’Abadia de Montserrat, 1997.}. В статье о каталонском крестьянском доме 1993~г. последний понимался как структура «окружения» и подчинения крестьян, связывавшая семейные хозяйства с юрисдикцией сеньора и видоизменявшая локальные общины\footnote{\emph{Idem.}~Le mas catalan du XIIe s.~: genèse et évolution d’une structure d’encadrement et d’asservissement de la paysannerie // Cahiers de Civilisation Médiévale. 1993. №~142 (36e année). P.~151--177.}. Наконец, в работе о крепостном праве и социальной мобильности Л.То~Фигерас показал, как динамика XI--XIII~вв. заложила социально-правовые предпосылки каталонских «дурных обычаев»\footnote{\emph{Idem.}~Servitude et mobilité paysanne~: les origines de la «Remença» catalane (Xie--XIIIe siècle) // Mélanges de l’école française de Rome. 2000. №~2 (112). P.~827--865.}.

С.~Собрекес-и-Видаль работал в каталонской историко-правовой и институциональной традиции, для которой характерно внимание к эволюции форм «производства права» и длительной преемственности от вестготского и римско-канонического права к каталонским обычаям и кодификациям. Именно ему принадлежит несколько обобщающих монографий, в которых «Барселонские обычаи» помещаются в широкую перспективу генезиса каталонских правовых сборников\footnote{\emph{Sobrequés Vidal S., Font i Rius J.M.~Història} de la producció del dret català fins al Decret de Nova Planta. Girona: Universitat de Girona, 1978; \emph{Sobrequés Vidal S., Peláez M.J.} Historia general del derecho catalán hasta el siglo XVIII. Barcelona: Promociones y Publicaciones Universitarias, 1989.}. Этот интерес к институциональной истории наметился в его более ранних исследованиях: о муниципальном режиме позднего Средневековья и системе выборов чиновников по жребию\footnote{\emph{Idem.}~Régimen municipal en la Baja Edad Media: La «insaculación» // Annals de l’Institut d’Estudis Gironins. 1955. №~10. P.~165--234.}. Эти изыскания позволили реконструировать механизмы функционирования городских советов и процедур политического представительства, а исследования по истории каталонской знати и «баронов» уточнили роль сеньориальных структур в отправлении правосудия\footnote{\emph{Idem.}~La nobleza catalana en el siglo XIV // Anuario de estudios medievales. 1970. №~7. P.~513--532.}. Исследования С.~Собрекеса-и-Видаля по истории средневековой Жироны дополнили эту проблематику анализом социальной структуры города и муниципального управления. Автор проследил изменения в организации городских органов власти, порядке их формирования и представительстве различных групп городского населения\footnote{\emph{Idem.}\emph{ }Societat i estructura política de la Girona medieval / ed. a cura de J. Sobrequés i Callicó. Barcelona: Curial, 1975.}.

Каталонский историк права Ж. Масип-и-Фонойосса является одним из ключевых исследователей обычного права Тортосы. Его многолетняя работа в Историческом архиве города сопровождалась систематическим изучением рукописной традиции «Обычаев Тортосы». Обнаружение ранее неизвестной рукописи XIV~в., исследование истории формирования памятника и подготовка его критического издания стали важнейшими результатами этой работы. В работах «Становление “Обычаев Тортосы”» и «Заселение Тортосы: предпосылки и исторический контекст» он проследил эволюцию правовых установлений этого городского сообщества в контексте формирования поселенческой структуры и соотношения королевской, сеньориальной и муниципальной юрисдикций\footnote{\emph{ Massip Fonollosa J.} La gestació de les costums de Tortosa. [Tortosa]: Consell Intercomarcal de les Terres de l’Ebre, 1984. 430 p.; \emph{Idem.} La Població de Tortosa: antecedents i context històric // Les Cartes de població cristiana i de seguretat de jueus i sarraïns de Tortosa (1148--1149): Actes de les Jornades d'Estudi commemoratives del 850è aniversari de la seva concessió, Tortosa, 14, 15 i 16 de maig de 1999. 2000. P.~151--168.}. Исследования «Органы отправления правосудия в Тортосе от поселенной хартии до обычаев» и «Администрация правосудия в Тортосе с 1149~г.» позволили проследить институциональную среду действия местного обычного права и механизмы его применения в судебной практике\footnote{\emph{Idem.} Els organs de la administració de justicia a Tortosa des de la carta de població a les costums // Jaime I y su época. 1979. T.~2. P.~461--473.}. В статье о «морских обычаях Тортосы» Дж.~Масип-Фонойосса показал, как торгово-морские общины адаптировали общегородские нормы к специфике средиземноморской коммерции. Источниковедческое измерение его исследований особенно заметно в анализе топонимики земель Эбро\footnote{\emph{Idem. }Problemes de la toponímia de les Terres de l’Ebre // Societat d’Onomàstica: butlletí interior. 1984. №~17. P.~9--11.} и характеристике рукописи «Обычаи Тортосы: критическое издание», найденной на Майорке\footnote{\emph{Idem.} Un manuscrit de Les Costums de Tortosa retrobat a Mallorca // Revista de llengua i dret. 1986. №~7. P.~61--68.}, что позволяет реконструировать географию действия и каналы передачи норм обычного права.

Особое место в историографии занимали работы А.~Гурона, который, обращаясь прежде всего к проблеме соотношения римского права и «общего права», фактически исследовал формирование средневекового европейского права. В статье «К истокам влияния глоссаторов в Испании» он проследил механизмы проникновения учёного права глоссаторов на Пиренейский полуостров и тем самым позволил увидеть, как учёное право начинало взаимодействовать с местными правовыми практиками\footnote{\emph{Gouron~A.} Aux origines de l' influence des glossateurs en Espagne // Historia. Instituciones. Documentos. 1983. №~10. P.~325--346.}. В другой работе «Являлось ли общее право наследником римского права?» он показал, что следует говорить о сложной интеллектуальной переработке римского материала в университетской среде, сопровождавшейся переосмыслением самих понятий «право», «обычай», «общее право»\footnote{\emph{Idem. }Le droit commun a-t-il été l’héritier du droit romain? // Les séances de l’Académie des Inscriptions et Belles-Lettres. 1998. Vol.~142. №~1. P.~283--292;}. Каталонские локальные обычаи выступали у А.~Гурона преимущественно как сопоставительный материал, позволявший выявить различия в темпах и формах рецепции римско-канонической традиции по обе стороны Пиренеев.

Т.Н.~Биссон работал в англо-американской традиции институциональной и социальной медиевистики, а также в постоянном диалоге с подходами школы «Анналов». Его статьи конца 1970-х гг. о «мире и перемирии Божием» на территории Каталонии и проблеме становления феодальной монархии задали политико-правовое направление его дальнейшим исследованиям\footnote{\emph{Bisson T.N.~The} Organized Peace in Southern France and Catalonia, ca. 1140 -- ca. 1233 // The American Historical Review. 1977. №~2 (82). P.~290--311; \emph{Idem. }The Problem of Feudal Monarchy: Aragon, Catalonia, and France // Speculum. 1978. №~3 (53). P.~460--478.}. В программной статье о «феодальной революции» Т.Н.~Биссон предложил хронологию и механизм трансформации властных отношений в XI--XII~вв., что позволило объяснить гибкость их социально-правовой терминологии\footnote{\emph{Idem.}~The «Feudal Revolution» // Past \& Present. 1994. №~142. P.~6--42.}.\textbf{ }В статье «Торжество и убеждение: размышления о культурной эволюции средневековой совещательной практики» Т.Н.~Биссон проследил изменение роли средневековых собраний: наряду с публичным выражением согласия с властью они постепенно становились местом обсуждения спорных вопросов и выражения различных интересов. На каталонском материале исследователь связывал этот процесс с политическими конфликтами конца XII~в.\footnote{\emph{Idem.}~Celebration and Persuasion: Reflections on the Cultural Evolution of Medieval Consultation // Legislative Studies Quarterly. 1982. Т.~VII. №~2 (May, 1982). P.~181--204.}. В монографии 1998~г. Т.Н.~Биссон совершил «антропологический поворот», отойдя от классической институциональной истории и рассмотрев сельские конфликты и повседневные судебные разбирательства как «голоса обывателей», через которые проявлялась кодификация и применение обычая\footnote{\emph{Idem.}~Tormented Voices: Power, Crisis, and Humanity in Rural Catalonia, 1140--1200. Cambridge; London: Harvard University Press, 1998.}. В синтетической концепции «кризиса XII~в.» этому исследователю удалось связать усиление власти графов Барселонских с возникновением общеевропейской концепции управления сельскими сообществами\footnote{\emph{Idem.} The Crisis of the Twelfth Century: Power, Lordship, and the Origins of European Government. Princeton; Oxford: Princeton University Press, 2015.}.

В совокупности и в хронологическом, и в тематическом разрезе исследования Т.Н.~Биссона показали, что эволюция каталонского законодательства объясняется не только автономным «созреванием» права, но и изменением геополитики, фискально-административного инструментария и совещательных практик. Тем самым для изучения обычного права XII--XIII~вв. «Барселонские обычаи» предстали как живой текст, в котором термины, юрисдикции и процедуры переопределяются под воздействием кризисов и консолидации власти графов-королей.

П.Х.~Фридман исследовал становление зависимого положения каталонского крестьянства и его правовое оформление. Его работы конца 1970-х --- начала 1980-х гг. заложили его интерпретацию происхождения рабства и механики сеньориальных прав на материалах монастырских и орденских архивов и сопоставления норм и практик\footnote{\emph{Freedman P.H.} An Unsuccessful Attempt at Urban Organization in Twelfth-Century Catalonia // Speculum. Vol.~54. №~3. 1979. P.~479--491; \emph{Idem.}~Military Orders and Peasant Servitude in Catalonia: Twelfth and Thirteenth Centuries // The Hispanic American Historical Review. 1985. Т.~65. No.~1. P.~91--110; \emph{Idem. }Peasant servitude in the thirteenth century // Estudi General. 1986. №~5--6. P.~437--445; \emph{Idem.} La condition des paysans dans un village catalan du XIIIe siècle // Annales du Midi. 1982. Т.~94. №~158. P.~231--244.}. Тематически их продолжает блок работ, в которых продемонстрировано, как доктрина перерабатывала «обычай» в легитимирующий нарратив\footnote{\emph{Idem.}~Catalan Lawyers and the Origins of Serfdom // Mediaeval Studies. 1986. Т.~48. №~1. P.~288--314; \emph{Idem. }The Catalan Ius Maletractandi // Recueil de mémoires et travaux publié par la Société d’Histoire du Droit et des Institutions des Anciens Pays du Droit Écrit. Montpellier : Faculté de droit et des sciences économiques de Montpellier, 1985. №~13. P.~39--53.}. Параллельно П.Х.~Фридман обратился к проблематике границ различных юрисдикций: исследование о пограничном статусе реки Льобрегат в XIII~в. показало, как пограничье структурирует компетенции власти и практики принуждения\footnote{\emph{Idem.}~The Llobregat as a frontier in the thirteenth century // Miscel·lània en homenatge al P.~Agustí~Altisent. Tarragona: Diputación Provincial de Tarragona, 1991. P.~109--118.}. В работах о легендарных истоках Каталонии П.Х.~Фридман показал, каким образом нарративы доблести и трусости встроены в легитимацию власти и обычаев\footnote{\emph{Idem.}~Cowardice, Heroism and the Legendary Origins of Catalonia // Past \& Present. 1988. Vol.~121. P.~3--28.}.

Эти направления исследования были теоретически обобщены в монографии «Происхождение крепостного права в Средневековой Каталонии»\footnote{\emph{Idem.} The Origins of Peasant Servitude in Medieval Catalonia. Cambridge; New York: Cambridge University Press, 1991.}. В этой работе П.Х.~Фридман задал хронологию и механизмы институционализации зависимости от XI к XIII~вв. на основе обширного комплекса источников (актовые, нарративные, правовые)\footnote{Подробнее об этой монографии на русском языке см.: \emph{Филиппов}\emph{ }\emph{И}\emph{.}\emph{С}\emph{. }Рец.: Freedman P. The Origins of Peasant Serfdom in Medieval Catalonia. Cambridge; New York, 1991 // СВ. 1993. Т.~56. С.~272--273.}. В поздних очерках он вернулся к теме кодификации права при Жауме I и к установлению «дурных обычаев» (в т.~ч. на материале Сан-Пере-де-Кассаррес), показав длительность существования и гибкость элементов обычного права\footnote{\emph{Freedman P.H.} L’elaboració del dret i el règim senyorial en el regnat de Jaume I // Jaume I: commemoració del VIII centenari del naixement de Jaume I (L’economia rural~; L’articulació urbana~; Les institucions eclesiàstiques~; L’expansió territorial~; El comerç). Barceona: Institut d’Estudis Catalans, 2013. Vol.~2. P.~61--71; \emph{Idem.}~ “Remences” and the Lordship of Sant Pere de Casserres // La Corona catalanoaragonesa, l’Islam i el món mediterrani: estudis d’història medieval en homenatge a la Doctora María Teresa Ferrer i Mallol. Barcelona : CSIC, Institució Milà i Fontanals, Departament de Ciències Històriques-Estudis Medievals, 2013. P.~277--282.}. В методологическом отношении для настоящего исследования важны точный терминологический анализ, строгая работа с датировками и последовательностями актов, а также умение связать локальные тяжбы с более широкой эволюцией правовой культуры.

Работы Д.Дж.~Кэгэя выдержаны в русле англо-американской традиции изучения права и институтов Арагонской Короны. Он рассмотрел каталонское прав сквозь призму политической истории, военной организации и становления арагоно-каталонской монархии. Во вступительной статье к английскому переводу «Барселонских обычаев»\footnote{\emph{Kagay D.J.} Introduction // The Usatges of Barcelona: The Fundamental Law of Catalonia. Philadelphia, 1994.~P.~3--66.} он показал, что компиляции римского права южнофранцузского и североитальянского происхождения задали «законодательный стиль» придворных легистов, стремившихся облечь нормы местного обычного и феодального права в традиционные римские формулы, а также выявил текстологические и стилистические аллюзии на «Вестготскую правду», подчеркнув неоднородность корпуса и необходимость его поэтапной, хронологически дифференцированной реконструкции\footnote{Ibid. P.~18.}. Фактически эта вступительная статья представляет собой самостоятельное монографическое исследование по истории каталонского права XI--XII~вв., в котором автор не только обобщил историографические дискуссии о датировке «первоначального ядра» и позднейших слоёв «Обычаев», но и связал хронологию их формирования с изменением баланса сил между графской властью и феодальной знатью\footnote{Ibid. P.~3.}. В последующих работах он проследил, как нормы, близкие по типу к «Барселонским обычаям», функционируют в городской и военной политике: это исследования королевской власти в городах и формирования административного аппарата и войска (статьи о Жауме I и городах\footnote{\emph{Idem.}~Royal Power in an Urban Setting: James I and the Towns of the Crown of Aragon // Mediaevistik. 1995. Vol.~8. P.~161--170; \emph{Idem.} The line between memoir and history: James I of Aragon and the Llibre del Feyts // Mediterranean Historical Review. 1996. P.~165--176.}, о военной организации Арагонской Короны и управлении\footnote{\emph{Idem.}~Army Mobilization, Royal Administration, and the Realm in the Thirteenth-Century Crown of Aragon Brill, 1996. P.~95--115; \emph{Kagay D.J.} War, government, and society in the medieval Crown of Aragon. Aldershot [England] ; Burlington, VT : Ashgate/Variorum, 2007.}).

Другой тематический блок составляют работы о королевском правосудии и политических конфликтах, где Д.Дж.~Кэгэй проанализировал судебные процессы, правовые формулы и королевские привилегии как инструменты переговоров между королём и политическими сообществами, тем самым выстраивая связь раннесредневекового обычая и формировавшейся структуры принципата Каталония\footnote{\emph{Idem.}~Rebellion on trial: The Aragonese Union and its Uneasy connection to royal law, 1265--1301 // The Journal of Legal History. 1997. Vol.~18. P.~30--43; \emph{Idem.}~The Treason of Center and Periphery: The Uncertain Contest of Government and Individual in the Medieval Crown of Aragon // Mediterranean Studies. 2003. Vol.~12. P.~17--35.}.

А.Дж. Косто в ряде работ рассматрел «Барселонские обычаи» в связи с грамотами-соглашениями, клятвами и иными способами письменного оформления обязательств между графами Барселонскими и знатью в XI--XIII~вв.\textbf{ }Анализ грамот типа «соглашение» (лат. --- convenientia) позволил А.Дж.~Косто проследить формулы, санкции и практики взаимных гарантий, через которые выстраивались отношения графов Барселонских и знати\footnote{\emph{Kosto A.J.} The"convenientiae" of the Catalan counts in the eleventh century: a diplomatic and historical analysis //Acta historica et archaeologica mediaevalia. 1998. №~19. P.~191--228.}.

В статье «Ограниченный характер действия “Барселонских обычаев” в Каталонии XII~в.» он показал слабую раннюю цитируемость свода, его элитарную направленность и медленное внедрение в практику --- в том числе из-за программно-политического характера текста, в отличие от имевшего более широкое хождение перевода «Книги приговоров» на старокаталанский\footnote{\emph{Idem.} The limited impact of the «Usatges de Barcelona» in Twelfth-century Catalonia // Traditio. Cambridge: Cambridge University Press, 2001. Т.~56. P.~53--88.}. Эта аргументация была расширена исследователем в монографии «Как договаривались в средневековой Каталонии: власть, порядок и письменность, 1000--1200~гг.», в которой центральное место отдано письменной фиксации норм обычного права при составлении грамот\footnote{\emph{Idem.} Making agreements in medieval Catalonia: power, order and the written word, 1000--2000. Cambridge: Cambridge University Press, 2001. P.~14.}.

В работе о «Большой книге феодов» (лат. --- \emph{Liber}\emph{ }\emph{feudorum}\emph{ }\emph{maior}) автор представил этот известный картулярий как средство легитимации притязаний графов Барселонских. В другой статье об участии мирян и клириков в составлении грамот А.Дж.~Косто проследил связь социальной структуры средневековой Каталонии и процессов записи и фиксации \footnote{\emph{Idem.} Laymen, clerics, and documentary practices in the early Middle Ages: the example of Catalonia // Speculum. 2005. Т.~80. №. 1. P.~44--74.}.

Сравнительное исследование практик составления документов в Каталонии и Англии около 975--1050~гг. поместило каталонский опыт в общеевропейский контекст «обращения к документу», проиллюстрировав специфическое сочетание графской и монастырской письменности и иные процедуры доверия к актам по сравнению с англосаксонской средой. В главе в коллективной монографии 2024~г. он подчеркнул, что одни и те же лица могли выступать и свидетелями, и поручителями, и арбитрами, и соприсяжниками, что объясняет устойчивость локального обычая и медленное внедрение графского или позднее королевского права через кодифицированные тексты\footnote{\emph{Idem.} Versatile Participants in Medieval Judicial Processes: Catalonia, 900--1100 // Records and Processes of Dispute Settlement in Early Medieval Societies: Iberia and Beyond. 2024. P.~311--337.}. Эти наблюдения подкрепили вывод о том, что после завоеваний в Новой Каталонии графам (а именно Рамону Беренгеру IV) требовались способы регулирования отношений со знатью и что «Барселонские обычаи» выступили скорее декларацией графской программы. При этом А.Дж.~Косто пришёл к закономерному умозаключению: регулирование сеньориальных отношений осуществлялось, прежде всего, посредством заключения локальных соглашений.

Представитель компаративистики Т.~де~Монтагут-и-Эстрагес сочетал институциональную историю с тщательной источниковедческой работой. Его обобщающие труды по истории каталонского права и учреждений аккумулировали достижения предыдущих десятилетий и задали устойчивую оптику историко-правовых исследований\footnote{См. обобщающие работы, задуманные одновременно и как монографии, и как справочно-обучающие материалы для студентов-историков и правоведов: \emph{Montagut}\emph{ }\emph{i}\emph{ }\emph{Estragués}\emph{ T. de, }\emph{Maluquer}\emph{ de Motes }\emph{i}\emph{ Bernet C.J. }Història del dret espanyol. Barcelona: Edicions de la Universitat Oberta de Catalunya, 1997.; \emph{Montagut i Estragués T. de, Ferro Delgado V., Serrano Daura J.} Història del dret català. Universitat Oberta de Catalunya, 2009.}. Он проследил механизмы включения каталонских обычаев в систему «общего права»: от исследования отчуждения феода без согласия сеньора (1211--1330~гг.)\footnote{\emph{Montagut i Estragués T. de.} La recepción del derecho feudal común en Cataluña I ( 1211--1330): (La alienación del feudo sin el consentimiento del Señor) // Glossae: European Journal of Legal History. 1992. №~4. P.~9--145.} и использования «Семи партид» Альфонсо X Мудрого в каталонском праве\footnote{\emph{Idem.} Presencia de Las Partidas en el derecho catalán // Espacios y fueros en Castilla-La Mancha (siglos XI--XV): una perspectiva metodológica. 1995. №~2. P.~487--505.} до общеевропейской постановки вопроса\footnote{\emph{Idem.} La cultura jurídica europea del ius commune (S. XII--XV): la teoría y la práctica de sus juristas // Ius Publicum. 2008. №~20. P.~39--51.} и поиска элементов  «общего права» в локальных установлениях г. Тортосы\footnote{\emph{Idem.} La Carta de Tortosa i el dret comú // Les Cartes de població cristiana i de seguretat de jueus i sarraïns de Tortosa (1148-1149): Actes de les Jornades d'Estudi commemoratives del 850è aniversari de la seva concessió, Tortosa, 14, 15 i 16 de maig de 1999, Barcelona: Universitat Internacional de Catalunya, 2000. P.~169--177.; \emph{Idem.} La recepción del derecho feudal común en Cataluña I ( 1211--1330 ): ( La alienación del feudo sin el consentimiento del Señor ) // Glossae: European Journal of Legal History. 1992. №~4. P.~9--145.}. Параллельно он реконструировал роль локальных обычаев в социальной истории барселонского общества XI--XIII~вв.\footnote{\emph{Idem.} La societat de Barcelona i el seu dret: Segles XI--XIII // Catalan Historical Review. 2008. №~1. P.~183--193.}

К источниковедческому направлению можно отнести исследования об эволюции правовых норм в средневековой Барселоне\footnote{\emph{Idem.} Recognoverunt proceres Barchinone, et antiqui et sapientes in iure Diputación Provincial de Tarragona, 2020. P.~121--144;\emph{ Idem.} Els usages de Barcelona, un text jurídic medieval de doble transició // La justícia a la Catalunya medieval: actes del VII Seminari d'Estudis Medievals d'Hostalric : (18 i 19 de novembre de 2021). Hostalric: Ajuntament d'Hostalric, 2022. P.~85--93.}, в ходе которых была уточнена датировка ключевых норм городского права столицы графства и продемонстрирован плавный переход от обычного к «ученому» праву. Тем самым Т.~де~Монтагут-и-Эстрагес предложил модель совместной эволюции обычного права и ученой доктрины. Наконец, в исследованиях об образе мусульманина в праве и литературе\footnote{\emph{Kagay D.J.} The Essential Enemy: The Image of the Muslim as Adversary and Vassal in the Law and Literature of the Medieval Crown of Aragon / \emph{D.R.~Blanks}, M. Frassetto, New York: Palgrave Macmillan US, 1999. P.~119--136.}, о Пере Альберте как теоретике феодального права\footnote{\emph{Idem. }Pere Albert: Barcelona Canon, Royal Advocate, Feudal Theorist // Anuario de Estudios Medievales. 2002. №~1 (32). P.~39--74.} и о громком уголовном процессе XIV~в.\footnote{\emph{Idem.}\emph{ }The Murder of the Abbot: A Homicide and its Wider Impact in Fourteenth-Century Catalonia // Mediaevistik. 2008. Vol.~21. P.~87--111.} он показал, как язык права и судебная практика стали ключевыми ресурсами формирования политической идентичности и закрепления новых представлений о преступлении, верности и государственном суверенитете. Эти работы вписали «Барселонские обычаи» и другие локальные своды обычного права в широкую хронологическую и сравнительную рамку англо-американской историографии, трактующей их не как изолированный памятник, а как один из этапов в длительной эволюции правовой культуры и публичного права Каталонии от XI до XIV~вв.

А.~Иглесиа Феррейрос --- один из представителей каталонской историко-правовой традиции изучения средневекового права, берущей начало от Ф.~Вальс-и-Табернера, Дж.М.~Фонт-Риуса и др. Для неё характерны внимание к эволюции правовых институтов Каталонского принципата и тщательный историко-филологический анализ правовых текстов.

В качестве руководителя междисциплинарной группы в Автономном университете Барселоны он на протяжении многих лет координировал проект подготовки критического издания «Барселонских обычаев», результатом которого стали публикации текста по большому числу рукописей с развёрнутыми палеографическими и юридическими комментариями, хотя полностью завершённого критического издания до сих пор не существует\footnote{\emph{Iglesia Ferreirós A.} Glosas y Usatges: el ms. BNP Latin 4670 A Edición // Initium: Revista catalana d’historia del dret. 2002. №~7. P.~363--847; \emph{Idem.} La compilación catalana del «ms». E // Initium: Revista catalana d’historia del dret. 2005. №~10. P.~381--544; \emph{Idem. }Variaciones sobre un mismo tema: los usatges de Barcelona // Initium: Revista catalana d’historia del dret. 2006. №~11. P.~3--105 и др. его работы.}. Ранняя статья исследователя «Право в раннесредневековой Каталонии»\footnote{\emph{Idem.} El derecho en la Cataluña altomedieval // Memorias de la Real Academia de Buenas Letras de Barcelona. 1991. Vol 24. P.~27--34.} задала широкую тематическую и методологическую рамку --- корпус источников, институциональный контекст, взаимодействие «высокого» (римско-канонического, вестготского) и локального права. Центральное место в его творчестве занимают исследования по происхождению и структуре «Барселонских обычаев» («Маленькая загадка --- происхождение “Барселонских обычаев”»\footnote{\emph{I}\emph{dem.}Un pequeño enigma el origen de los “Usatici” // Mundos medievales: espacios, sociedades y poder~: homenaje al profesor José Ángel García de Cortázar y Ruiz de Aguirre. 2012. Vol.~1. P.~521--609.}, «Рождение “Обычаев”»\footnote{\emph{Idem.} The Birth of the Usatici // Imago temporis. Medium Aevum. 2011. №~5. P.~119--134.} и др.). В этих работах он последовательно выделил различные редакционные слои корпуса, сопоставив их с ключевыми этапами политической истории графства Барселонского XII--XIII~вв. и пересмотрев традиционные датировки раннего ядра и последующих добавлений\footnote{\emph{Idem.} Giraud, D’Abadal y Valls, Mor y los Usatges // Initium: Revista catalana d’historia del dret. 2002. №~7. P.~3--78.}. В этих работах А.~Иглесиа Феррейрос сопоставил нормы «Обычаев» с данными других каталонских и арагонских правовых сборников, сравнил манускрипты между собой, уточнил время фиксации отдельных титулов и предложил варианты решения дискуссионных вопросов об авторстве и первоначальном назначении этого правового памятника.

Отдельный блок его публикаций посвящён месту «Барселонских обычаев» в правовой системе Каталонии и Арагонской Короны, взаимодействию обычного, феодального и королевского права, а также рецепции римско-канонических формул в тексте памятника. Как историк права, он последовательно сочетает юридический анализ норм с классическими источниковедческими и компаративными методами, что позволяет ему вписать «Барселонские обычаи» в более широкий контекст формирования правовой культуры средневековой Каталонии. Результаты его исследований представляют собой важный этап осмысления спорных источниковедческих и содержательных проблем и служат методологической основой для современного историко-правового изучения как непосредственно «Барселонских обычаев», так и каталонского обычного права XII--XIII~вв. в целом.

Каталонский медиевист Ф.~Сабатэ-и-Куруль работает в направлении франко-испанской традиции социальной и политико-правовой истории, сочетающей институциональный анализ с вниманием к языку источников и антропологии права. Несмотря на то, что большинство его работ посвящено позднесредневековому этапу развития каталонского общества, уже в статье «Структурные основы городской олигархии Каталонии» он заложил основы для изучения городской элиты и механизмов её обособления\footnote{\emph{Sabaté i Curull F.} Ejes vertebradores de la oligarquía urbana en Cataluña // Revista d’Historia Medieval. 1998. №~9. P.~127--149.}. В работе 1999~г. Ф.~Сабатэ-и-Куруль обратился к эволюции институтов управления, прослеживая формирование пространств власти и компетенций, в которых затем функционировало и обычное право\footnote{\emph{Idem.} La Governació al Principat de Catalunya i als comtats de Roselló i Cerdanya // Anales de la Universidad de Alicante. 1999. №~12. P.~21--62.}. Его программная статья по историографии представляла собой подробный обзор теорий и моделей становления феодального общественного устройства в трёх основных исторических регионах Испании (в том числе Каталонии)\footnote{\emph{Idem.}  L'apparition du féodalisme dans la péninsule Ibérique. Etat de la recherche au commencement du XXIe siècle // Cahiers de Civilisation Médiévale. 2006. Vol.~49, Nº 193. P.~49--69.}. Обращаясь к «Барселонским обычаям» и другим памятникам обычного права, автор одновременно подвергает критическому анализу концепции своих предшественников и систематизирует предложенные ими объяснения региональных различий. В исследовании «Смертная казнь в позднесредневековой Каталонии: портрет общества» он уже на уровне судебной практики проследил соотношение писаных норм и местных обычаев, показав, каким образом репрессивные механизмы государства и города переосмысляли традиционные представления о справедливости и наказании\footnote{\emph{Idem.} La pena de muerte en la Cataluña bajomedieval // Clio \& Crimen. 2007. №~4. P.~118--276.}. Ф.~Сабатэ-и-Курулем неоднократно обращался к лексике актового и нарративного материала, фиксируя, как сами каталонцы обозначали властные полномочия и формы зависимости.

Дж.М.~Сальрак-и-Марес как представитель каталонской социальной и политико-правовой истории, во многом связанной с французской традицией «Анналов», также занимался исследованием средневекового каталонского крестьянства. В ранних работах о крестьянстве и сеньориальной власти он рассмотрел эволюцию зависимых состояний от IX до XIII~вв., тем самым задав социально-исторический фон для анализа обычного права этого периода\footnote{\emph{Salrach i Marès J.M.} Remences i pagesos a la Catalunya medieval // Revista d’historia medieval. 1995. №~6. P.~127--135.}. Статья 2001~г. дополнила этот блок, связав изменения в аграрном обществе и развитии рынков с эволюцией правовой коммуникации и механизмов защиты прав\footnote{\emph{Idem.} Sociedad rural y mercados en la Cataluña medieval // Edad Media: revista de historia. 2001. №~4. P.~83--111.}. Особое место среди его работ занимают исследования судебной практики и разрешения конфликтов: на материале процедурных форм, компромиссов и арбитража Дж.М.~Сальрак-и-Марес проследил постепенное оформление политических инстанций и разграничение их компетенций\footnote{\emph{Idem.} Prácticas judiciales, transformación social y acción política en Cataluña (siglos IX-XIII) // Hispania: Revista española de historia. 1997. №~197 (57). P.~1009--1048; \emph{Idem.}  Les modalités du règlement des conflits en Catalogne aux XIe et XIIe siècles // Actes des congrès de la Société des historiens médiévistes de l’enseignement supérieur public. 2000. №~31. P.~117--134.}.

В поздней монографии «Правосудие и власть в Каталонии до 1000~г.» исследователь сместил акцент на источниковедческий анализ, реконструируя ранние (IX--XI~вв.) представления о справедливости, на основе которых затем сложилась правовая система XII--XIII~вв. В этой работе особенно заметно внимание автора к языку актов, что делает её методологически значимой для настоящего исследования. Наконец, в статье «Государственное и феодальное правосудие в Каталонии» Дж.М.~Сальрак-и-Марес подвёл итоги размышлений о многослойности судебной системы и предложил понимаычного права между сеньориальными, королевскими и коммунальными практиками\footnote{\emph{Salrach i Marès J.M.} Justícia d’Estat i justícia feudal a Catalunya // Butlletí de la Societat Catalana d’Estudis Històrics. 2022. №~33. P.~279--302.}.

В историографии обычного права Новой Каталонии особое место занимает творчество Дж.~Серрано Дауры, который последовательно рассмотрел локальное обычное право пограничных территорий (Тортоса, Льейда и др.) как ключ к пониманию механизмов их интеграции в каталонское политико-правовое пространство. В ряде работ о распространении права Льейды, о сеньориальных и муниципальных структурах Орты и Миравета, а также об «Обычаях Фликса» и «Обычаях Тортосы» он показал, как сформировались локальные своды обычаев, сочетающие привнесённые модели с практиками повседневного регулирования на новом заселённом пространстве. Исследования Дж.~Серрано Дауры построены на тщательном источниковедческом анализе конкретных законодательных памятников --- от поселенных хартий до полноценных сводов обычного права, что позволяет реконструировать не только текстуальную историю норм, но и контекст их применения. Особый интерес для изучения локального обычая представляют его работы о мультиконфессиональной среде и социальных и религиозных практиках морисков. В этих исследованиях локальное обычное право выступило инструментом упорядочения сосуществования различных общин. Анализируя соотношение сеньориальной власти и муниципального самоуправления, а также участие знати в отправлении правосудия, Дж.~Серрано Даура показал, что через судебные процедуры, четкое определение границ юрисдикций и активное участие городских корпораций в составлении и записи локальных сводов обычного права проявились особенности обычного права Новой Каталонии\footnote{\emph{Serrano Daura J.} Usos i pràctiques socials i religiosos dels moriscos d’Ascó, Benissanet i Miravet // Miscel·lània del Centre d’Estudis de la Ribera d’Ebre. 1993. №~9 (1). P.~31--42; \emph{Idem.} L’antic dret de Lleida a les comarques tarragonines (s. XII-XIV) 1997. P.~533--540; \emph{Idem. }Senyoria i municipi a Orta: s. XII-XIII 1997. P.~85--101; \emph{Idem.} Incorporació definitiva a Catalunya dels territoris d’Ascó i de Miravet (1347), i d’Horta (1359) 1999. P.~67--95;\emph{ Idem.} La Carta de Seguretat deis sarraïns de Tortosa, de 1148 2000. P.~105--150; \emph{Idem.} Els costums de Flix, una baronia de Barcelona // Barcelona: quaderns d’història. 2001. №~5. P.~87--102; \emph{Idem. }Senyoria i municipi a la Batllia de Miravet (s. XII-XIV) 2002. P.~97--137; \emph{Idem.} La projecció del dret de Lleida // Revista de dret històric Català. 2004. (3). P.~233--258; \emph{Idem. }El dret marítim en els Costums de Tortosa (1277--1279) València: Universidad de Valencia = Universitat de València, 2005. P.~569--582; \emph{Idem. }La famille dans l’historiographie juridique des territoires hispaniques pyrénéens 2005. P.~51--78; \emph{Idem.} La Coexistència de les comunitats cristiana, jueva i sarraïna a Tortosa a la baixa edat mitjana // Revista de Dret Històric Català. 2006. №~6. P.~173--193; \emph{Idem. }Le droit catalan au temps de la Croisade contre les Albigeois 2010. P.~99--112; \emph{Idem. }El Judici de prohoms, una institución judicial de participación vecinal // Glossae. 2015 (12). P.~782--800;\emph{ Idem.} La concessió dels costums de la batllia de Miravet // Revista de dret històric Català. 2015 (14). P.~273--304; \emph{Idem. }Una aproximació als Costums de Tortosa i el seu contingut (1277--1279) // Revista de Dret Històric Català. 2023. №~22. P.~111--154.}.

Работы португальского исследователя Р.Р.~Тостеса можно отнести к классической пиренейской школе истории права, для которой характерно внимание к понятийному аппарату, институциональному языку и длительным хронологическим срезам. Его исследования посвящены широкому кругу вопросов: искусству и технике права в аргументации необходимости народного представительства при Пере IV Церемонном (1336--1387)\footnote{\emph{Tostes R.R.} Els artificis i la tècnica del dret dins dels arguments de la representativitat en el regnat de Pere el Cerimoniós // Revista de dret històric Català. 2018. Vol.~17. P.~81--117.}, институциональному языку каталонских кортесов XIV~в.\footnote{\emph{Idem.} Autoridad, comunidad política y representación: los cambios semánticos y una mirada hacia atrás en la Cataluña medieval // El parlamentarisme en perspectiva històrica. Parlaments multinivell. 2019. Vol.~2. P.~1027--1047.}, эволюции понятий «авторитет», «политическое сообщество»\footnote{\emph{Idem.} Una lectura sobre el lenguaje institucional en las asambleas parlamentarias catalanas del siglo XIV // Calamus. 2019. Vol.~3. P.~102--119.}, фигуре судьи-монарха\footnote{\emph{Idem.} El monarca-juez y las greuges en las cortes catalanas de la Edad Media // Cortes y Parlamentos en la Edad Media peninsular. 2020. P.~473--500.} и др.

Особое место среди них занимает работа, в которой автору удалось связать внутреннюю хронологическую структуру локальных сводов обычного права с изменениями правовой культуры\footnote{\emph{Idem.} Do seny i bon juhi na criacao do direito publico: a mutacao do costume na catalunha medieval (seculos XII--XIII) // Signum. Revista da ABREM. 2020. Vol.~21. №~1. P.~86--112.}.

Важное значение для настоящего исследования имеет статья каталонских историков права Ф.~Удина-и-Марторелю и А.М.~Удина-и-Абельо «Соображения относительно “первоначального ядра” “Барселонских обычаев”»\footnote{\emph{Udina i Martorell~F., Udina i Abelló~A.M.} Consideracions a l’entorn del nucli originari dels “Usatici Barchinonae” // Estudi general: Revista de la Facultat de Lletres de la Universitat de Girona. 1985--1986. №~5--6. P.~87--107.}, в которой эта проблема рассматривается не только в плоскости датировки, но и как ключ к пониманию механизмов формирования правовой культуры Каталонии.

Проследив эволюцию теории «первоначального ядра» с момента её возникновения до 1970-х гг., авторы показали, как менялись представления исследователей о соотношении письменной фиксации обычая, судебной практики и графской власти. Их источниковедческий подход, основанный на детальном анализе рукописной традиции, формулировок отдельных обычаев и аргументации средневековых юристов, позволил реконструировать не только хронологические слои текста, но и представления современников о природе «обычая» и его авторитетности. В полемике с ранними исследователями они настояли на отказе от механического поиска «наиболее древних» норм и подчеркнули необходимость учитывать динамику переработки и переосмысления обычного права в учёной среде.

Особое место занимает работа каталонского исследователя А.~Масферрера-Доминго, касающаяся влияния «Барселонских обычаев» на уголовно-правовое регулирование в муниципиях Новой Каталонии\footnote{\emph{Masferrer Domingo A.} La influència dels Usatges en el Dret penal dels municipis de la Catalunya Nova (notes per a un estudi) // El territori i les seves institucions històriques. Actes. Ascó, 28, 29 i 30 de novembre de 1997. Barcelona. 1999. P.~809--838.}. Он использовал историко-правовой подход, анализируя рецепции и адаптации «Барселонских обычаях» в городском праве Льейды, Орты, Миравета и других городских общин. Рассмотрение истории формирования текста и историографической традиции, сложившейся вокруг него, позволило вписать каталонское городское право XIII--XIV~вв. в более широкий контекст эволюции обычного права Пиренейского полуострова.

Другой каталонский специалист по истории права В.~Гарсия-Эдо последовательно рассмотрел каталонские компиляции обычного права не изолированно, а в широком контексте формирования правовых систем Арагонской короны в целом и Королевства Валенсии в частности. В ряде работ он показал, как валенсийские сборники XIII~в. возникали и развивались в тесной связи с политикой Жауме I и распространением королевского права Арагонской короны\footnote{\emph{García Edo V}. Origen i expansió dels Furs o Costum de València, durant el regnat de Jaume I // Boletín de la Sociedad Castellonense de Cultura. 1993. №~69. P.~175--200; \emph{Idem.} La redacción y promulgación de la «Costum» de Valencia // Anuario de estudios medievales. 1996. №~26. P.~713--728; \emph{Idem.} Territorialización del derecho foral valenciano // Revista valenciana d’estudis autonòmics. 2008. №~51. P.~104--143.}. Сравнительный анализ обычаев Льейды, Валенсии и Тортосы позволил В.~Гарсии-Эдо выявить текстологические связи соответствующих сборников обычаев и проследить трансформацию каталонских правовых моделей при их адаптации в валенсийской среде\footnote{\emph{Idem. }El parentesco entre las costumbres de Lérida (1228), Valencia (1238) y Tortosa (1273) // Anuario de historia del derecho español. 1997. №~67. P.~173--188.}. В исследованиях об источниковой базе «Обычаев Тортосы» и влиянии на них итальянского права, связанного с полномочиями нотариев, он связал эволюцию местного обычного права с рецепцией римского права\footnote{\emph{Idem.} Fonts documentals dels Costums de Tortosa i la influència de la notorialística italiana del dret comú a la rúbrica «De notaris» // Boletín de la Sociedad Castellonense de Cultura. 2000. T.~76. №~1--4. P.~37--56.}. Обращаясь к латинской версии «Фуэро Хаки», В.~Гарсия-Эдо поместил каталонские и валенсийские компиляции в более широкий арагонский контекст\footnote{\emph{Idem.} La versión latina del Fuero Extenso de Jaca del siglo XII (Una propuesta de reconstrucción del texto) // Aragón en la Edad Media. 2017. №~28. P.~39--66.}. В результате каталонские записи обычного права он понимал как ключевое, но не единственное звено в развитии средневекового права Арагонской Короны, сформировавшейся на пересечении каталонской, арагонской и валенсийской традиций.

Американский исследователь П.~Дэйладер, работающий в русле англо-американской традиции социальной и политической истории, внёс существенный вклад в понимание механизмов функционирования правовой культуры Каталонии в период расширения территорий и консолидации власти графов-королей. В своей программной статье он показал, как локальные своды обычного права, возникавшие в различных городах и регионах от Руссильона до Тортосы, вступали в диалог с «Барселонскими обычаями», одновременно адаптируя, дополняя и переосмысляя их нормы в соответствии с конкретными социальными и конфессиональными условиями\footnote{\emph{Daileader Ph.} La coutume dans un pays aux trois religions~: la Catalogne, 1228 -- 1319 // Annales du Midi~: revue archéologique, historique et philologique de la France méridionale. 2006. №~225 (118). P.~369--385.}.

П.~Дэйладер продемонстрировал, что процесс кодификации обычного права был плюралистическим: записи местных обычаев интегрировали элементы из «Вестготской правды», римского права и канонического права, создавая многослойную социальную структуру, отражавшую сложность управления мультиконфессиональным обществом. Особенно значимо его наблюдение о том, что некоторые городские компиляции (особенно в Руссильоне) формировались в полемике с «Барселонскими обычаями», что свидетельствовало о динамичном характере правовой коммуникации и отсутствии жёсткой иерархии между центральным и локальным правом\footnote{\emph{Idem.} Town and countryside in northeastern Catalonia, 1267 -- c. 1450: the sobreposats de la horta of Perpignan // Journal of Medieval History. 1998. №~4 (24). P.~347--366.}. Его исследования, основанные на анализе документов из городских архивов и законодательных сводов, позволили рассмотреть распространение и трансформацию обычного права как процесс постепенного переосмысления управленческих практик и институциональных структур\footnote{\emph{Idem.} True Citizens: Violence, Memory, and Identity in the Medieval Community of Perpignan, 1162--1397. Leiden; Boston: Brill, 2000.}.

В XX---начале XXI~вв. появляется много исторических исследований, авторы которых пользуются данными «Барселонских обычаев» в качестве основного или дополнительного источника в рамках узких проблемных и тематических областей. Среди них -- история женщин и семьи (работы М.~Ауреля-Кардоны\footnote{\emph{Aurell Cardona M.} Du nouveau sur les comtesses catalanes (IXe--XIIe siècles) // Annales du Midi. Vol.~109. №~219--220.}, М.~Комас-Виа\footnote{\emph{Comas-Via M.} Looking for a Way to Survive: Community and Institutional Assistance to Widows in Medieval Barcelona // Women and Community in Medieval and Early Modern Iberia. Lincoln: University of Nebraska Press, 2020. P.~177--194.}), история взаимодействия различных религиозных сообществ в средние века (Е.~Кляйн\footnote{\emph{Klein E.} Les juifs et la société chrétienne à Barcelone au Moyen Âge // Annales du Midi. 2006. Vol.~118. №~225. P.~437--440.}, П.~Сенак\footnote{\emph{Senac P.} Islam et chrétienté dans l’Espagne du haut Moyen Age: la naissance d’une frontière // Studia Islamica. 1999. №~89 (1999). P.~91--108.}, Г.~Гонсалво-и-Боу\footnote{\emph{Gonzalvo i Bou G.} Els jueus i els Usatges de Barcelona // Barcelona Quaderns D'Història. Barcelona, 1996. №~2. P.~117--124.}, А.~Гарсия-и-Гарсия\footnote{\emph{García y García~A.} Los juramentos e imprecaciones en los “Usatges” de Barcelona // Glossae: European Journal of Legal History. Instituto de Estudios Sociales, Políticos y Jurídicos, 1995. №~7. P.~51--80.}, Х.Х.~Мунтане-и-Сантивери\footnote{\emph{Muntané i Santiveri J.-X.} Anàlisi de l’estructura del jurament de les malediccions dels jueus catalans: usatge 171 // Revista de Dret Històric Català. Barcelona, 2014. Vol.~13. P.~9--48.} и др.), история судопроизводства и его особенности (Л.~Ассье-Андрьё\footnote{\emph{Assier-Andrieu~L.} Dificultad y necesidad de la antropología del derecho // Revista de Antropología Social, 1970. Vol.~24. P.~35--52; \emph{Idem}. Custom and Law in the Social Order: Some reflections upon French-Catalan peasant communities // Law and History Review. 1983. Vol.~1. №~1. P.~86--94.}, Л.~Донат Рохас\footnote{\emph{Donat Rojas L.} The duel in medieval western mentality // Ideology in the Middle Ages: Approaches from Southwestern Europe / ed. Sabaté i Curull. Leeds: Arc Humanities Press, 2019. P.~175--202.}, Ж.-М.~Рубеллан\footnote{\emph{Rubellin M. }Combattant de Dieu ou combattant du Diable? Le combattant dans les duels judiciaires aux IXe et Xe siècles // Actes des congrès de la Société des historiens médiévistes de l’enseignement supérieur public. Montpellier: Publications de la Sorbonne, 1987. Vol.~18. P.~111--120.} и др.).

Классические работы заложили основы научных представлениях о происхождении, содержании и структуре текстов каталонского обычного права. Именно к этому этапу относятся первые успешные попытки построения периодизаций его развития. При этом, большинство исследователей привлекали тексты записей обычного права как источник для конкретно-проблемных исследований, посвящённых различным вопросам, прежде всего, социально-экономической. Значительное влияние оказали историографические традиции такие, как например, школа «Анналов», определившие концептуально новые постановки вопросов и интерпретации источников.

\subsubsection*{§1.3. Терминологические исследования}
\addcontentsline{toc}{subsubsection}{§1.3. Терминологические исследования}

В настоящий подпараграф вошли работы по лексике отдельных сводов обычного права, актового материала, а также обобщающие терминологические исследования о бытовании важных для нашей диссертации понятий в контексте средневековой истории Пиренейского полуострова.

Особое значение для нашего исследования имеет монография каталонской исследовательницы Е.~Родон-Бинуэ «Специализированная лексика каталонского феодализма XI~в.»\footnote{\emph{Rodon Binue~E.} El lenguaje técnico del feudalismo en el siglo XI en Cataluña. (Contribución al estudio del latin medieval). Barcelona: Escuela de Filología, 1957. 278 p.}. В центре её внимания находилось формирование «технического» языка феодальных отношений, что позволило проследить, каким образом локальные обычаи и преимущественно устная практика заключения соглашений закреплялись в латинской документации. На основе обширного корпуса актового материала автор выявила формы написания ключевых терминов, их этимологию, частотность и контексты употребления, показав процесс постепенной фиксации повседневных правовых практик и их адаптации к потребностям нотариата и суда. Для настоящей диссертации особенно важно, что в работе Е.~Родон-Бинуэ наглядно выражен «пограничный» между разговорной и книжной традицией характер ряда понятий. Вместе с тем для прояснения происхождения отдельных терминов, связанных с социальной структурой каталонского общества, оказалось необходимым обратиться к исследованиям, посвящённым аналогичной терминологии в других регионах (в первую очередь в Северной Италии\footnote{\emph{Rippe G.} Dans le Padouan des Xe--XIe siècles~: évêques, vavasseurs, «cives~» // Actes des congrès de la Société des historiens médiévistes de l’enseignement supérieur public. Poitiers: Publications de la Sorbonne, 1983. Vol.~14. P.~141--150.}).

В том же русле, обозначающем поворот к истории правового дискурса и правовой ментальности, выполнены работы К.~Дуарте-и-Монсеррата --- прежде всего исследование «Правовая лексика “Обычаев Тортосы” (MS. de 1272)», а также статьи о языке и диалектных особенностях данной рукописи\footnote{\emph{Duarte i Montserrat C.} El vocabulari jurídic del Llibre de les costums de Tortosa (ms. de 1272). 1a ed. Barcelona: Generalitat de Catalunya, Departament de Governació, Escola d’Administració Pública de Catalunya, 1985. 128 p.; \emph{Idem.} La llengua del manuscrit de 1272 del Llibre de les Costums de Tortosa // Revista de llengua i dret. 1991. №~16. P.~7--56; \emph{Idem.} Notes sobre trets dialectals en el «Llibre de les costums de Tortosa» // Caplletra. Revista Internacional de Filologia. 1987. №~2. P.~19--26.}. Посвящённые конкретному компендиуму --- «Обычаям Тортосы», они показали, каким образом при заимствовании норм и формул из авторитетных сборников (например, «Барселонских обычаев») происходила их языковая и концептуальная адаптация к местной среде. К.~Дуарте-и-Монсеррат проследил, какие термины, обозначающие институты, процедуры и статусы, перешли в состав текста «Обычаев Тортосы», какие подверглись переопределению или уточнению и как соотносились латинская и старокаталанская традиции.

Французская исследовательница А.~Рюкуа вместе с обобщающими работами по истории Пиренейского полуострова выступила автором главы в коллективной монографии под редакцией Ф.~Сабатэ-и-Куруля «Auctoritas, potestas: концепты власти в средневековой Испании»\footnote{\emph{Rucquoi~A. }Auctoritas, potestas: concepts of power in medieval Spain // Ideology in the Middle Ages: Approaches from Southwestern Europe / ed. Sabaté i Curull~F. Leeds, 2019. P.~51--72.}. Этот раздел важен не столько кастильским материалом, сколько обращением к истории понятий. А.~Рюкуа показала, как латинские “auctoritas” и “potestas” в христианских королевствах Пиренейского полуострова семантически меняли свои значения и встраивались в систему представлений о законности, иерархии и легитимации власти. Проанализировав королевские акты, нормативные сборники и нормативные тексты, она продемонстрировала, как римско-правовые категории интегрировались в практику и повседневные представления о праве.

Диссертация каталонского филолога К.А.~Прието-Эспиносы «Профессиональная лексика в средневековой каталонской документации на латинском языке»\footnote{\emph{Prieto Espinosa C.A.} El lèxic dels oficis a la documentació llatina de la Catalunya altmedieval. / Tesi presentada per optar al Grau de Doctor. Barcelona. 2020. 2 vol.} вывела изучение текстов права за пределы сугубо лексического анализа и вписала в более широкую линию исследований каталонского обычного права, обозначенную, в частности, работой Е.~Родон-Бинуэ.

Используя методы корпусной лингвистики и опираясь на «Словарь средневековой каталонской латыни», построенный как корпус текстов (включая онлайн-публикацию «Барселонских обычаев» --- 2020), К.А.~Прието-Эспиноса рассмотрел правовую лексику как индикатор складывающейся профессиональной среды, институциональных структур и практик взаимодействия между различными группами общества. Профессиональные термины, обозначения должностных лиц, участников судебных процедур, видов договоров и форм насилия он понял не только как языковые единицы, но и как носителей определённых правовых представлений и социальных иерархий. В этом смысле его глоссарий и наблюдения над функционированием отдельных слов и словоформ оказываются ценными как вспомогательный инструмент при работе с латинской и старокаталанской версиями «Барселонских обычаев» и других сборников локального обычного права.

Тем самым лингвистический анализ понятийного аппарата «Обычаев Тортосы» превращается в исследование отдельных правовых норм и понятий, позволяя увидеть, как в городской среде XII--XIII~вв. сформировалась собственная терминология, опирающаяся на авторитет общекаталонских сборников обычного права, таких, как «Барселонские обычаи», но перерабатывающая их в соответствии с локальными практиками. К.А.~Прието-Эспиноса и К.Дуарте-и-Монсерат рассматривали нормы обычного права в связи с формой их записи, употреблением терминов и особенностями их функционирования в тексте. Тем самым проблематика, ранее представленная в исследованиях Е. Родон-Бинус, получила дальнейшее развитие.

К началу 2000-х гг. исследовательское внимание сконцентрировалось на письменных практиках, что позволило уточнить механизмы записи обычного права. Сложившаяся панорама начала XXI~в. демонстрирует продуктивный союз институциональной истории, социально-экономической и правовой истории для объяснения изменений, произошедших в XII--XIII~вв.

Историография каталонского обычного права формировалась на пересечении двух крупных исследовательских традиций --- французской школы, методологически близкой «Анналам», и англо-американской институционально-социальной медиевистики. Представители первой (П.~Боннасси, Ж.М.~Сальрак-и-Марес, отчасти Ф.~Сабате-и-Куруль) рассматривали локальные своды обычного права как выражение длительных процессов феодализации, трансформаций сеньориальной власти и развития социальных структур с акцентом на анализе актового материала. Для англо-американских исследователей (Т.Н.~Биссон, А.Дж.~Косто, П.~Фридман) обычное право предстало прежде всего результатом политико-административных практик, переговоров и документальной культуры, что позволяло переосмыслить «Барселонские обычаи» и другие записи обычного права как инструмент власти, а не универсальные модели кодификации. Эти две школы задали разные, но взаимодополняющие исследовательские оптики, определив хронологические рамки анализа каталонского обычного права XII--XIII~вв.

Сравнение французской и англо-американской историографических традиций особенно важно для исследования каталонского обычного права, поскольку оно позволяет увидеть, как одни и те же источники можно трактовать посредством различных методологических подходов. Французская школа подчеркнула социальную природу локального обычая и его место в длительных процессах феодализации, тогда как англо-американская традиция выявляет механизмы производства норм через письменность, административные практики и политические коммуникации. Сопоставление этих подходов помогает избежать одностороннего объяснения и выработать комплексное понимание того, как формировалась и функционировала правовая культура Каталонии XII--XIII~вв. Тем самым сравнительный анализ становится методологическим основанием для исследования каталонских правовых текстов, включая «Барселонские обычаи».

Терминологические исследования показали, что язык описания средневекового права не является нейтральным и существенно влияет на интерпретацию источников. Эти разработки создают важную теоретическую и методологическую основу для анализа источников, используемых в диссертации.

Обзор западноевропейской медиевистики демонстрирует, что интерес к каталонскому обычному праву прошёл путь от фрагментарных упоминаний к комплексным исследовательским программам. Постепенно фокус сместился от реконструкции правовых институтов к изучению практик, дискурсов и культурных кодов. Значимую роль сыграли работы, интегрирующие методы правовой истории, социальной истории и исторической антропологии.

\subsection*{§2. Социолингвистический подход к изучению обычного права}
\addcontentsline{toc}{subsection}{§2. Социолингвистический подход к изучению обычного права}

В западноевропейской медиевистике во второй половине XX~в. возрос интерес к многоязычным письменным практикам. Активно изучается, каким образом латинский язык взаимодействовал с народными в правовых текстах --- как повседневных, так и нормативных. Перевод и языковые изменения в исследовательском сообществе стали пониматься как социально-культурный процесс, а не как сугубо техническая задача---присвоение новых правовых смыслов и дискуссия о статусе нормы. В этой части историографического раздела рассматриваются исследования мультилингвизма и языковой ситуации в Средиземноморье и на Пиренейском полуострове, а также работы, посвящённые соотношению устного и письменного в средневековом праве. Особое место занимают исследования средневекового перевода и интерпретации.

\subsubsection*{§2.1. Методологические основания изучения мультилингвизма}
\addcontentsline{toc}{subsubsection}{§2.1. Методологические основания изучения мультилингвизма}

Базовые понятия теории мультилингвизма задали классические тексты, посвящённые диглоссии и билингвизму. Ч.A.~Фергюсон в программной статье «Диглоссия» описал устойчивое распределение функций между «высоким» и «низким» кодами. Собственно диглоссию он определил как явление, подразумевающее сосуществование двух языков с разделением сфер их использования (в нашем случае старокаталанского языка и латыни, использовавшегося в официальных документах). Ч.A.~Фергюсон отметил, что такая языковая ситуация может иметь разное происхождение и проявляться в разных языковых ситуациях. Кроме того, он увидел перспективность применения этого понятия по отношению к исторической социолингвистике\footnote{\emph{Ferguson C.A.} Diglossia // WORD. Journal of the International Linguistic Association Routledge, 1959. Vol.~15. №~2. P.~326.}.

Немного позже Дж.~Фишман разделил «билингвизм» как индивидуальную компетенцию (ведущую своё происхождение от психолингвистических исследований) и «диглоссию» как социальную организацию языковых кодов, которая представляла интерес для исследователей этого периода\footnote{\emph{Fishman J.A. }Bilingualism with and without Diglossia; Diglossia with and without Bilingualism // Journal of Social Issues. 1967. №~2 (23). P.~29--30.}.

Для анализа микромеханизмов переключения кодов применяется концепция интерактивной социолингвистики Дж.Дж.~Гамперца \emph{“}\emph{we-code}\emph{/}\emph{they-code}\emph{”}. Эта теория иллюстрирует положение о том, что переключение кодов является сложным социальным явлением в зависимости от сферы деятельности, адресата и характера высказывания\footnote{\emph{Gumperz}\emph{ J.J.} Discourse strategies. Cambridge; New York~: Cambridge University Press, 1982. P.~143.}.

Отличиям в использовании заимствований и переключения кодов посвящена статья лингвистов С.~Поплака и Д.~Санкофа «Заимствование: одновременность интеграции»\footnote{\emph{Poplack}\emph{~S., Sankoff~D.} Borrowing: the synchrony of integration. // Linguistics. Vol.~22. Jena, 1984. P.~101.}. Это исследование позволило провести границы между заимствованиями отдельных лексических единиц из старокаталанского языка и включением этих единиц в синтаксические структуры и грамматические рамки латинского предложения.

К 1980-м гг. сформировалось несколько вариантов подходов к сугубо лингвистическим исследованиям билингвизма и переключения кодов. Однако исследовательские интересы лингвистического сообщества в 1990-х гг. сменили направление и предложили взгляд на переключение кодов, смешение нескольких языков и другие сходные явления в контексте формирования идентичности отдельных социальных и этнических групп в динамике\footnote{\emph{Myers-Scotton C. }Duelling languages: Grammatical structure in code-switching. Oxford: Oxford Universiry Press, 1993. P.~2; \emph{Auer P.} From codeswitching via language mixing to fused lects // International Journal of Bilingualism. 1999. Т.~3. №~4. P.~1--28.}.

В конце XX--- начале XXI~в. вектор научных интересов исследователей сместился от случаев диглоссии к случаям многоязычия --- мульти- и плюрилингвизма. Среди концептуальных работ, посвящённых этому явлению, стоит указать на главу «Переключение кодов и грамматическая теория» (П.~Миускен)\footnote{\emph{Muysken}\emph{ P.} Code-switching and grammatical theory // One Speaker, Two Languages: Cross-Disciplinary Perspectives on Code-Switching / ed. Milroy L., Muysken P.~Cambridge, 1995. P.~85.} в коллективной монографии, представляющей собой междисциплинарное исследование явления мультилингвизма. В данной работе автор приводит формальные модели, согласно которым строятся предложения в разговорной практике людей, владеющих двумя и более языками. Важным для нашего исследования наблюдением П.~Миускена является тот факт, что в сознании говорящего всегда есть «доминирующий» и «подчинённый» языки и что наиболее часто мультилингвизм проявляется в использовании отдельных лексических и синтаксических единиц «доминирующего» языка внутри предложения на подчинённом языке.

В коллективной монографии, изданной в 2000~г. под редакцией П.В.~Кроскрити, был введён термин «режим языка», который описывал идеологические конструкции, придающие общему (например, в нашем случае латинскому языку) трансрегиональную легитимность, а вернакулярам --- локальную аутентичность и социальную адресность\footnote{\emph{Kroskrity}\emph{ P.V.} Regimes of Language: Ideologies, Polities, and Identities. d. Santa Fe: School of American Research Press, 2000.}.

Дж.Н.~Адамс в работе «Билингвизм и латинский язык»\footnote{\emph{Adams J.N.} Bilingualism and the Latin Language. Cambridge University Press, 2003.} коснулся многих принципиальных для нашего исследования проблем, а именно соотношения латинского языка и вернакуляра в текстах, имеющих различный социальный контекст и синтаксической обусловленности латинских вставок и заимствований в текстах на иных языках, исторических принципов переключения кодов для различных групп языков, в том числе и иберо-романских.

\subsubsection*{§2.2. Мультилингвизм в средневековом Средиземноморье и на Пиренейском полуострове}
\addcontentsline{toc}{subsubsection}{§2.2. Мультилингвизм в средневековом Средиземноморье и на Пиренейском полуострове}

Обратимся к исследованиям, где указанный инструментарий применяется к средневековому и пиренейскому материалу. Ключевым шагом для преодоления традиционной дихотомии «высокого» и «низкого» регистров использования латыни и вернакуляра стала «социофилология» Р.~Райта\footnote{\emph{Wright R}. Sociolinguistique hispanique (VIIIè -- XIè siècles) // Médiévales. 1993. Т.~12, №~25. P.~69;\emph{ Idem. }Latin and the Romance languages in the early Middle Ages. University Park: Pennsylvania State University Press, 1996; \emph{Idem.} Translation between Latin and Romance in the Early Middle Ages // Translation Theory and Practice in the Middle Ages. / ed. \emph{Beer~J.M.}A.~Medieval Institute Publications, Western Michigan University, 1997. P.~7--32.}. Он систематизировал идею континуума между позднелатинскими и раннероманскими практиками письма, в том числе с использованием пиренейского материала (преимущественно не каталонского). В своих исследованиях Р.~Райт проанализировал специфику позднелатинской письменности различных областей Пиренейского полуострова. Автор подчеркнул, что пути складывания местных диалектов во многом определялись политическими факторами. Например, сравнив позднелатинские тексты на территории Кастилии и Каталонии, он сделал вывод, что в последнем случае довольно длительная зависимость от государства Каролингов в качестве Испанской марки способствовала сближению с провансальскими и окситанскими вариантами латинского языка в Каталонии, тогда как кастильская разновидность испытала значительное арабское влияние\footnote{Ibid.}.

При этом Р. Райт определил языковую ситуацию в Каталонии к X~в. как «диглоссию»\footnote{Ibid. P.~68--69.}. В документах «латынь» часто выступает как совокупность жанрово обусловленных регистров, часть из которых фактически опирается на живую романскую грамматику и лексику. Р.~Райту удалось показать это на примере текстов различных жанров, в которых была распространена формульность, где «нормативная» латинская грамматика соседствовала с романскими морфемами и топонимикой. Такой взгляд радикально изменил исследовательский подход к «вульгаризмам» в актовом материале и нормативных текстах: это не ошибки, а показатели социолингвистического диапазона писца и общества в целом.

Весьма полезным для нас стало исследование лексики актового материла на вульгарной латыни XI~в. авторства Ж.-М.~Тока\footnote{\emph{Tock B.-M.} Les mutations du vocabulaire latin des chartes au XIe siècle // Bibliothèque de l’École des chartes. Paris, 1997. Vol.~155. №~1. P.~119--148.}. Анализируя документальный корпус, автор выявляет наиболее частотные лексические изменения, характерные для интересующего нас периода. Несмотря на то что автор проводит сквозной поиск по леммам в корпусе текстов, он подчёркивает необходимость тщательной работы с каждым текстом и документом, отмечает важность внимания к этимологии слов, их семантическому окружению, синтаксису и т.д.\footnote{Ibid. P.~146.} Работа выполнена на южно-французском материале, однако в данном случае нам важно сочетание двух подходов --- сплошного поиска в массиве тысяч текстов определённых лексических элементов и скрупулёзная аналитическая работа автора (пример case-studies), связанная с обращением к историческому и языковому контексту составления, бытования и функционирования актового или делопроизводственного материала.

К проблеме перехода от народной латыни к романским языкам обращаются Ж. Бельмон и Ф. Вейярд в работе «“Фаршированная латынь” и окситанский язык в актах XI~в.»\footnote{\emph{Belmon J., Vielliard F.} Latin farci et occitan dans les actes du XIe siècle // Bibliothèque de l’École des chartes. Paris, 1997. Vol.~155. №~1. P.~149--183.}. Авторы считают, что на начальном этапе данной трансформации единицы последнего включались в документы, составленные на вульгарной латыни, посредством использования местной лексики, вписанной в латинские грамматические и синтаксические конструкции. Это и есть «фаршированная латынь». Только когда актовый материал и документы, связанные с судопроизводством, насыщаются этой лексикой, происходит формирование литературного, но не регулярного народного романского языка, однако на данном этапе на первый план уже выходят региональные особенности.

Соотношению устного и письменного в средневековых правовых практиках посвящена статья каталонского историка Х.В.~Боски-и-Кодины «От~голоса к тексту: изменение и постоянство в процессе письменной фиксации (Каталония, X--XII~вв.)»\footnote{\emph{Boscá Codina~J.V.} De la voz en el texto: cambios y permanencias en el proceso de afirmación de la escritura (Cataluña, siglo X--XII) //Acta historica et archaeologica mediaevalia. 1999. №~20. P.~139--175.}. В данной работе автор анализирует то, каким образом устные формулы попадали в нормативные тексты, обращается он и к «Барселонским обычаям» как к одному из вспомогательных источников.

Исследованием историко-филологического характера является монография К.Е.~Леглу «Мультилингвизм и родной язык в средневековых французских, окситанских и каталанских нарративных источниках»\footnote{\emph{Leglu}\emph{ C.E}. Multilingualism and Mother Tongue in Medieval French, Occitan and Catalan Narratives. Pennsylvania; Penn State University Press, 2010.}. Несмотря на то что круг источников данного труда практически не связан с юриспруденцией и её особенностями, в этой работе содержатся многие концептуальные наблюдения, касающиеся сосуществования нескольких языков в интересующем нас регионе, особенностей графики и записи текстов на «долитературном» этапе формирования местных романских языков, которые оказываются очень полезными в контексте изучения существовавших в средневековой Каталонии переводческих практик.

Мультилингвизм в Средние века остается одной из наиболее востребованных и интересующих исследователей различных регионов (к нашему сожалению, практически не затрагивая Пиренейский полуостров). Среди новейших публикаций, так или иначе посвящённых применению лингвистической оптики мультилингвизма, стоит назвать сборник «Языки раннесредневековых грамот (латынь, германские вернакуляры и писаное слово)»\footnote{The Languages of Early Medieval Charters. Latin, German Vernaculars, and the Written Word / ed. R. Gallagher, E. Roberts, F. Tinti. Boston; Leiden: Brill, 2021. 548 p. (Brill’s Series on the Early Middle Ages, T.~27)}: многие статьи, вошедшие в его состав, демонстрируют методологическую преемственность с классическими работами Р. Райта, М.Т.~Клэнчи и др.\footnote{\emph{Tinti F.} Latin and Germanic Vernaculars in Early Medieval Documentary Cultures: Towards a Multidisciplinary Comparative Approach // The Languages of Early Medieval Charters. Latin, German Vernaculars, and the Written Word / Ed. by R. Gallagher, E. Roberts, F. Tinti. Boston; Leiden: Brill, 2021. P.~3.}

В другом сборнике «Языки и сообщества в западных провинциях поздней Римской империи и в постимперский период»\footnote{Languages and Communities in the Late-Roman and Post-Imperial Western Provinces / Ed. by Mullen A., Woudhuysen G. Oxford: Oxford University Press, 2023.} представлены две статьи, посвящённые анализу языковой ситуации на Пиренейском полуострове в раннем Средневековье, написанные в том числе с применением методик изучения многоязычных сообществ, посредством анализа текстов и отдельных языковых единиц (также на основе данных ономастики)\footnote{\emph{Velázquez I.} Reflections on the Latin Language Spoken and Written in Visigothic Hispania // Languages and Communities in the Late-Roman and Post-Imperial Western Provinces / ed. by Mullen A., Woudhuysen G. Oxford: Oxford University Press, 2023. P.~58--84; \emph{Barrett G.} Conservatism in Language: Framing Latin in Late Antique and Early Medieval Iberia // Ibid. P.~85--125.}.

Таким образом, методология исследования письменных практик многоязычных обществ уже оформилась, была успешно апробирована на средиземноморском и пиренейском материале и при этом остаётся востребованной в современной медиевистике. В совокупности эти исследования формируют представление о многоязычной среде Средиземноморья и Пиренеев, что позволяет нам перейти к анализу взаимодействия устной и письменной традиций в правовых практиках.

\subsubsection*{§2.3. Соотношение устной и письменной традиции в средневековых правовых текстах}
\addcontentsline{toc}{subsubsection}{§2.3. Соотношение устной и письменной традиции в средневековых правовых текстах}

В современной западной медиевистике проблема соотношения устного и письменного в источниках права рассматривается как комплексное явление, охватывающее не только (и не столько) технику фиксации нормы, но и социальные механизмы её функционирования. Исследования последних десятилетий (М.~Клэнчи, Р.~Маккитерик, М.~Баньяр, М.~Зиммерман, У.~Браун, Б. Сток, П. Гири и др.) продемонстрировали, что письменная фиксация обычного права не вытесняла устных практик, а трансформировала их, создавая смешанные формы правовой коммуникации.

Представление о тесной связи письменной фиксации с устными формами правового действия сложилось в английской историографии задолго до работ М.Т.~Клэнчи. Уже Ф.~Поллок и Ф.У.~Мейтленд обращали внимание на значение устного волеизъявления, свидетельства и публично совершаемых действий при передаче и удостоверении права\footnote{\emph{ Pollock F., Maitland F.W. }The History of English Law before the Time of Edward I. 2nd ed. Cambridge: Cambridge University Press, 1898. Vol.~2. P.~83--86.}. В середине XX~в. Ф.М.~Стентон рассматривал ранние англосаксонские грамоты в связи с практикой устных пожалований и сопровождавших их символических действий\footnote{\emph{Stenton F.M. }Anglo-Saxon England. 2nd ed. Oxford: Clarendon Press, 1947. P.~141}. П.~Шапле впоследствии существенно усложнил эту модель, поставив под сомнение представление о последовательной замене устного акта письменным документом\footnote{\emph{Chaplais}\emph{ P. }The Origin and Authenticity of the Royal Anglo-Saxon Diploma // Journal of the Society of Archivists. 1965. Vol.~3. No.~2. P.~49--52.}. М.Т.~Клэнчи, опираясь на сложившуюся исследовательскую традицию, обобщил и развил эти наблюдения на материале Англии XI--XIII~вв., проследив длительное сосуществование и взаимодействие устных и письменных форм правовой практики --- присяг, свидетельских показаний, публичных объявлений, символических действий и их документальной фиксации\footnote{\emph{Clanchy}\emph{ M.T.} Remembering the Past and the Good Old Law // History. 1970. Vol.~55. No.~184. P.~166--172; \emph{Idem.} From Memory to Written Record: England, 1066--1307\emph{.} Cambridge, Mass.: Harvard University Press, 1979. P.~116--117, 231--232. Также см. перевод статьи \emph{М.Т.~Клэнчи} на русский: \emph{Клэнчи}\emph{ }\emph{М}\emph{.} Воспоминания о прошлом и старый добрый закон // Vox medii aevi / под ред. \emph{Горбун Г.С.}; пер. \emph{Горбун Г.С.}, \emph{Заплатина А.А.}~СПб.: Перепечкин Кирилл Владимирович, 2015. №~2--3 (13--14). С.~112--128.}. Этот подход релевантен и для Каталонии XII--XIII~вв., где акты, сборники обычаев и муниципальные постановления были продуктом взаимодействия устного ритуала и письменной фиксации.

Б.~Сток предложил концепцию «текстуальных обществ», в рамках которой на примере вальденсов показал, что рост грамотности сопровождался сохранением устных форм правовой коммуникации. Важные тексты озвучивались в присутствии неграмотных, превращая чтение в публичное действо, такое как оглашение хартии, принесение присяги или публичное заседание королевского или графского суда. Всё это представляло собой гибридные устно-письменные практики, в которых текст обретал авторитет и легитимность в момент звучания\footnote{\emph{Stock B. }The implications of literacy: written language and models of interpretation in the eleventh and twelfth centuries. Princeton: Princeton University Press, 1983. P.~13--14.}.

В фундированной монографии «Живая речь: письменная и устная коммуникация на латинском западе с IV по IX век» М.~Баньяр рассмотрел оглашение, диалогичность и мнемотехнику как элементы, оставившие след в структуре и формулярах текстов актового материала поздней Античности и раннего Средневековья\footnote{\emph{Banniard M.} Viva voce: communication écrite et communication orale du IVe au IXe siècle en occident latin. Paris: Institut des études augustiniennes, 1992. P.~46--47. Эта работа была недавно переведена на итальянский язык и издана на нём с авторскими дополнениями. Подробнее см. \emph{Аникьев}\emph{ И.И., Филиппов И.С. }Banniard M. Viva voce: comunicazione scritta e comunicazione orale nell’occidente latino dal IV al IX secolo / Edizione italiana con una Retractatio dell’autore, a cura di L. Cristante e F. Romanini, con la collaborzione di J. Gesiot e V. Veronesi. Trieste: Edizione Università di Trieste, 2020. XVIII, 716 p. (Polymnia: Studi di filologia classica; 25). Баньяр М. Живая речь: письменная и устная коммуникация на латинском западе с IV по IX век / Итал. изд. с доп. автора; под ред. Л. Кристиани и Ф. Романини, при уч. Дж. Джезиота и В. Веронези. Триест: Издательство Университета Триеста, 2020. XVIII, 716 с. (Полимния: Исследования по классической филологии; 25) // СВ. 2021. Т.~82, №~3. С.~184--196.}. Он реконструировал «акустическую» сторону учёной латинской культуры, в которой чтение вслух, диктовка и повторение формул конструировали синтаксис, ритмику и грамматические структуры документов. Показав устойчивость этих практик, М.~Баньяр объяснил, почему формулы актов сохраняют следы публичного чтения. Стоит также отметить полемику М.~Баньяра и Р.~Райта. Р.~Райт критиковал первого за невнимание к корреляции между орфографией и конкретными фонетическими особенностями народных языков\footnote{\emph{Wright R.} Review Article: Michel Banniard, Viva Voce // The Journal of Medieval Latin. Turnhout: Brepols, 1993. №~3. P.~90.}.

Исследуя монастырские картулярии, П.~Гири выявил механизм «избирательного конструирования памяти» --- целенаправленной компиляции и отбора документов, где устные предания, локальные практики и их запись объединялись в рамках легитимирующей традиции\footnote{\emph{Geary P.J.} Phantoms of Remembrance: Memory and Oblivion at the End of the First Millennium. Princeton: Princeton University Press, 1996. P.~7.}. Показав, что картулярии представляли собой масштабные исторко-правовые проекты, он проследил стратегии переупорядочивания, сокращений и вставок, на службе у имущественных и юрисдикционных споров.

В процессе разработки проблемы дихотомии устного и письменного в каролингскую эпоху Р.Д.~Маккитерик выделила институциональную практику публичного чтения правовых актов, при которой документ функционировал как часть устной традиции\footnote{\emph{McKitterick R.D.} Introduction // The Uses of Literacy in Early Mediaeval Europe. Cambridge: Cambridge University Press, 1990. P.~3.}. Подчеркнув связку писаного слова с властью, она показала, что канцелярии и локальные суды создавали тексты с расчетом на их публичное воспроизведение --- от королевских капитуляриев до частных актов. В ее интерпретации публичное чтение не только распространяло нормы, но и придавало им силу через коллективное восприятие и запоминание. Этот вывод важен для понимания каталонской судебной сцены: оглашение, подтверждение общиной и материализация в записи выступают взаимодополняющими стадиями единого процесса, тем более что в течение нескольких веков Каталония находилась в сфере влияния каролингской документальной культуры.

М.~Зиммерман, анализируя каталонские актовый материал и текст «Барселонских обычаев», показал феномен «документальной революции» --- появления огромного массива актовых материалов при сохранении ключевых устных механизмов санкционирования: оглашение, присутствие свидетелей, участие общины\footnote{\emph{Zimmerman M.M. }Écrire et lire en Catalogne (IXe--XIIe siècle). Madrid: Casa de Velázquez, 2003. Vol.~1. P.~2. (Bibliothèque de la Casa de Velázquez; vol. 23); \emph{Idem.} L’usage du droit wisigothique en Catalogne du IXe au XIIe siècle: Approches d'une signification culturelle // Mélanges de la Casa de Velázquez. 1973. Т.~9. P.~233--281; \emph{Idem}. L’usage du droit wisigothique en Catalogne du IXe au XIIe siècle: Approches d'une signification culturelle // Mélanges de la Casa de Velázquez. 1973. tome 9. P.~233--281; \emph{Idem.} La représentation de la noblesse dans la version primitive des Usatges de Barcelone (milieu du XIIe siècle) // Cahiers de linguistique et de civilisation hispaniques médiévales. N° 25. Paris, 2002. P.~25.}. Указывая на региональные различия, он проследил профессионализацию писцов и постепенное оформление нотариальной практики, а также стандартизацию формул и знаков удостоверения под влиянием трансрегиональных моделей. Исследователь связывает резкое увеличение количества грамот с перестройкой властных отношений и режимов доказательства --- от «живого» свидетельства к поиску доказательства в письменных источниках нормативного или актового характера.

В русле англо-американской исследовательской традиции, опирающейся на правовую антропологию и социальную историю (Т.Н.~Биссон, Р.~Райт и др.), находится работа Д.А.~Боумена «Сменяющиеся ориентиры: собственность, доказательства и земельные тяжбы в Каталонии рубежа X--XI веков»\footnote{\emph{Bowman J.A.} Shifting Landmarks: Property, Proof, and Dispute in Catalonia around the Year 1000 / \emph{J.A.~Bowman}, Ithaca: Cornell University Press, 2004. 279 p.}. Исследуя практику разрешения земельных споров, Д.А.~Боумен показал, как сосуществовали и влияли друг на друга локальные обычаи, ритуалы клятвы и свидетельства с одной стороны и всё более сложный корпус письменных хартий --- с другой. Такой подход сблизил его с каталонской традицией изучения «пограничной зоны» между устным и письменным правопорядком (Е.~Родон-Бинуэ и др.), но он сделал акцент не столько на терминологии, сколько на процессуальной стороне конфликтов и практиках доказательства. Д.А.~Боумен показал, что рост значения письменной документации не вытеснял устное право, а перестраивал его, придавая устным формам новую правовую функцию и вписывая их в рамки развивающейся «документальной культуры». Для настоящего исследования его выводы особенно важны как модель анализа длительного взаимодействия письменного и устного в правовой культуре Каталонии: формы взаимодействия около 1000~г. задали контекст для понимания трансформаций обычного права и судебной практики в XII--XIII~вв.

Во вступлении к сборнику статей «Средневековый судебный процесс: физические, устные и письменные практики в Средние века» М.~Мостерт отметил, что письменность в правовой практике сосуществовала с ритуалом и устной вербализацией и что важно анализировать не только документы, но и связь между речью, действием и текстом\footnote{\emph{Mostert M.} Introduction // Medieval legal process: Physical, spoken and written performance in the Middle Ages / ed. M. Mostert, \emph{P.S.~Barnwell}. Turnhout: Brepols, 2011. P.~4--5. (Utrecht Studies in Medieval Literacy; T.~22)}.

В рамках того же сборника П.С.~Барнвелл рассмотрел прямую речь, включённую в тексты варварских правд, как нормативные образцы речевых актов, предназначенные для конкретных процессуальных ситуаций, а не как дословные записи реально произнесённых высказываний. Письменная фиксация, таким образом, лишь стандартизировала устный «перформатив», задавая формулы и сценарии поведения. Согласно его выводам, «устное» в нормативных текстах выступает как моделированная, идеализированная перформативность, тогда как «письменное» оформляет и регулирует ее, что ограничивает возможности прямой реконструкции фактической устной практики\footnote{\emph{Barnwell P.S. }Action, speech and writing in early Frankish legal proceedings // Medieval legal process: Physical, spoken and written performance in the Middle Ages / ed. M. Mostert, \emph{P.S.~Barnwell}. Turnhout, 2011. P.~23--24.}.

Указал на выборочный характер записи У.~Браун. В своём исследовании он установил, что документы, фиксирующие исход конфликтов, представляют собой ретроспективную редакцию устных споров. Фактически это дало нам возможность предположить, что обычай постоянно «переписывался» в нужном ключе и в интересах отдельных лиц\footnote{\emph{Brown W.C}. The gesta municipalia and the public validation of documents in Frankish Europe // Documentary Culture and the Laity in the Early Middle Ages / ed. by \emph{W.C.~Brown}, M.~Costambeys, M.~Innes and \emph{A.J.~Kosto}. Cambridge: Cambridge University Press, 2012. P.~95.}.

Очевидно, запись обычного права в средневековом обществе не может рассматриваться без учёта устной составляющей. В средневековых «текстуальных обществах» многочисленные документы включались в устный контекст: легитимность акта подтверждалась публичным оглашением его содержания, присутствием свидетелей, печатью и, следовательно, закреплением в коллективной памяти. Тем самым «обычай» подлежал фиксации, т.е. записи, но продолжал существовать как устная перформативная практика.

\subsubsection*{§2.4. Роль переводов с латыни на вернакуляры в Средние века}
\addcontentsline{toc}{subsubsection}{§2.4. Роль переводов с латыни на вернакуляры в Средние века}

Особого внимания заслуживают и обобщающие исследования о культурном и лингвистическом значении перевода в Средние века, которые показывают, как эти практики работали в конкретных институциональных и языковых контекстах и почему их логика не сводится к привычным для современной теории перевода представлениям.

Р.~Копелэнд показала, что средневековый перевод представлял собой форму ученого чтения и толкования текста. Он подчинялся школьным и жанровым нормам, опирался на риторические приёмы переработки авторитетного текста и был тесно связан с представлением о «перенесении учёности» (лат. --- \emph{translatio}\emph{ }\emph{studii}).\footnote{Подробнее о значении и истории термина см. \emph{Жоно}\emph{ Э. }Translatio studii: перенесение учёности (одна из тем Жильсона) // Интеллектуальные традиции Античности и Средних веков (исследования и переводы) / Сост. и общ. ред., пер. \emph{М.С.~Петровой}. М.: Кругъ, 2010. C.~238--294.} и становлением вернакулярной словесности. В соответствии с идеями Р.~Копелэнд, риторико-герменевтическая традиция задала основу для теоретических рассуждений о переводе и статусе языков\footnote{\emph{Copeland R.} Introduction // Rhetoric, Hermeneutics, and Translation in the Middle Ages: Academic Traditions and Vernacular Texts. Cambridge: Cambridge University Press, 1991. P.~5.}.

К.~Буридан показал на материале французского Средневековья, что переводы выступали двигателем языковой эволюции: через калькирование и семантическое заимствование они насыщали лексику, способствовали грамматикализации и нормированию, а также формировали дискурс о «правильности» письменного французского в сопоставлении с латынью\footnote{\emph{Buridant~C.} Le rôle des traductions médiévales dans l'évolution de la langue française et la constitution de sa grammaire // Médiévales. 2003. №~45. P.~77. URL: \url{http://journals.openedition.org/medievales/637}; DOI : \url{https://doi.org/10.4000/medievales.637É}, (last access: 02.05.2026).}.

В контексте изучения средневековых переводов с латыни на романские языки работа Л.~Бадии, Дж.~Сантанака и А.~Солера позволила рассмотреть вернакуляризацию как изменение коммуникативной стратегии распространения знания\footnote{\emph{Badia L., }\emph{Santanach}\emph{ J., Soler A. }Ramon Llull as a Vernacular Writer: Communicating a New Kind of Knowledge. Woodbridge: Tamesis, 2016. P.~235.}. На материале наследия Рамон Льюль\footnote{\textbf{Ра}\textbf{мон}\textbf{ }\textbf{Л}\textbf{ьюль} (1232--1316?) --- каталонский религиозный деятель, писатель, поэт, создатель каталанского литературного языка (см. Луллий Раймунд // Большая российская энциклопедия: научно-образовательный портал.  URL: \url{https://bigenc.ru/c/lullii-raimond-c3a7be/?v=8844626} (Дата обращения: 03.11.2023).} авторы показали, что латинская и каталанская версии одного и того же текста включались в единую систему производства и распространения знания, ориентированную на разные аудитории и адаптированную под потребности каждой из них. Такой подход особенно важен для анализа тех памятников, где перевод на романсе сопровождался перестройкой структуры памятника и даже его содержания.

К кастильскому материалу обратилась М.~Кастильо~Льюч, показав, как кастильские переводы правовых текстов XIII~в. (от двора Альфонсо X до толедских переводческих мастерских) создавали новую писаную норму через управление вариативностью: переводчики и писцы вырабатывали стратегию взаимодействия между латинскими и арабскими письменными моделями, унифицировали орфографию и синтаксис, закрепляли терминологию\footnote{\emph{Castillo Lluch M. }Translación y variación lingüística en Castilla (siglo XIII): la lengua de las traducciones // Cahiers d'études hispaniques médiévales. 2005. №~28. P.~131--144.}.

Согласно выводам А.~Бурё, перевод в Средние века был неотделим от практик записи, рубрикации, составления глосс и компиляции норм. Все эти технические этапы «фиксации «технические» процедуры структурировали смыслы не меньше, чем собственно межъязыковой перенос, и задавали рамки авторства, жанра и авторитетности\footnote{\emph{Boureau A. }L’écriture comme prisme. Contraintes et créations des scribes médiévaux // L’Atelier du Centre de recherches historiques. Revue électronique du CRH. 2017. №~16 Bis. URL: \url{https://journals.openedition.org/acrh/7479} (last acess: 04.11.2025).}.

Наконец, К.~Хёйш представил обзор основных работ, посвящённых переводу с латыни на романсе на Пиренейском полуострове, и отказался анахроничного отождествления средневекового перевода с современными нормами «эквивалентности»\footnote{\emph{Heusch C.} Penser la traduction au Moyen Âge. Problèmes et perspectives // Cahiers d’études hispaniques médiévales. 2018. Т.~41. №~1. P.~19.}. Этот автор призвал учитывать спектр значений в конфессиональном, правовом и научном дискурсах и развивать компаративные подходы в рамках теории мультилингвизма.

В совокупности эти работы задали понятийную рамку для трактовки перевода как многоуровневой интеллектуальной технологии, сочетающей интерпретацию, адаптацию и нормотворчество и потому неизбежно расходящейся с нашим сегодняшним представлением о «теории и методологии» переводов

\subsubsection*{§2.5. Исследования переводов латинского на старокаталанский язык в Средние века}
\addcontentsline{toc}{subsubsection}{§2.5. Исследования переводов латинского на старокаталанский язык в Средние века}

Наиболее близка к проблематике нашего исследования работа испанского филолога Л.~Рубио Гарсии «Сравнение латинского и каталанского языка “Барселонских обычаев”»\footnote{\emph{Rubio García~L.} Comparación entre el texto latino y el catalán de los Usatges de Barcelona // Glossae: European Journal of Legal History. 1995. №~7. P.~33--50.}. Эта статья написана исключительно на материале «Барселонских обычаев» и представляет собой сопоставление латинского и каталанского текстов, взятых по билингвальному изданию Ф.~Вальса-и-Табернера\footnote{Ibid. P.~33.}. Эта работа Л.~Рубио Гарсии характеризуется крайней конспективностью, которая отвечает заявленной цели его исследования: \emph{«<…> выявить возможные аномалии, лакуны и несовпадения, которые могут существовать между обеими редакциями текста <…>»}\footnote{Ibid.}. В выводах автор остановился на констатации отдельных изменений в переводе и особенностей старокаталанского языка, не объяснив их происхождение и не поместив их в широкий исторический и социолингвистический контекст.

В том же тематическом номере журнала «Глоссы: Европейский журнал правовой истории» была опубликована статья исследователей-филологов П.~Диеса~де~Ревенги и И.~Гарсии~Диаса, посвящённая сравнению различных вариантов перевода «Барселонских обычаев» на старокаталанский язык\footnote{\emph{Diez de Revenga~P., García Diaz~I.} Manuscritos y copistas: Los Usatici Barcinonae // Glossae: European Journal of Legal History. Murcia, 1995. №~7. P.~82.}. В данном исследовании учёные сопоставили способы перевода ряда статей «Барселонских обычаев» в разных вариантах, включая и взятый из рукописи, использованной Х.~Бастардесом-и-Парерой при издании «Барселонских обычаев». Авторы не обделили вниманием и лексические особенности старокаталанских переводов, попытались установить взаимосвязи среди латинских переводов памятника.

Испанский филолог Х.~Моран-и-Осеринхауреги\footnote{\emph{Morán i Ocerinjauregui~J}. El proceso de creación del catalán escrito // Aemilianense: revista internacional sobre la génesis y los orígenes históricos de las lenguas romances. San Millán de La Cogolla, 2004. Vol.~1. P.~437.} назвал «Барселонские обычаи» первым памятником «регулярного» старокаталанского языка, поскольку распространение и постоянное использование вернакуляра привели к их переводу в XIII~в. Исследователь установил, что центров переводов «Барселонских обычаев» могло быть несколько, на основе разночтений в старокаталанских рукописях памятника. Автор вписал появление перевода «Барселонских обычаев» на старокаталанский язык в общеисторический контекст эпохи. При этом Х.~Моран-и-Осеринхауреги указал на расширение территорий, на которых было распространён старокаталанский язык, связав это обстоятельство с экспансией Арагонской Короны в XIII~в. В приложении автор привёл фрагменты из старокаталанской версии «Барселонских обычаев» с собственными комментариями. Не менее важна и датировка первых переводов этого правового компендиума на старокаталанский XIII~в., поскольку она является одной из основных источниковедческих проблем, связанных с бытованием текста «Барселонских обычаев».

К числу современных лингвистических исследований «Барселонских обычаев» мы отнесли работы немецкого лингвиста Р.~Мейера-Германна\footnote{\emph{Meyer-Hermann R. }Auxiliar-Konstruktionen im llatímedieval katalanischer Urkunden (10.--13.~Jh.) und in den Usatici Barchinonae / Usatges de Barcelona (12. -- 15. Jh.) // Zeitschrift für romanische Philologie (ZrP). 2019. Bd. 2 (135). S. 507--534; \emph{Idem}. Zur Emergenz katalanischer Syntax in den Usatici Barchinonae / Usatges de Barcelona (12.--15. Jhdt.) // Zeitschrift für romanische Philologie (ZrP). 2020. Bd. 33. S. 155--199.}, в которых он описал ряд экспериментов, проведённых им на материалах текстов актов из Каталонии XI--XIV~вв. и рукописей «Барселонских обычаев» с применением технологий компьютерной лингвистики. В основном его работы посвящены синтаксическому переходу от вульгарной латыни старокаталанскому языку. Автор зафиксировал синтаксические особенности текстов, установил частоты изменения позиции подлежащего и сказуемого в рамках предложений. Р.~Мейер-Германн также обратился сразу к нескольким вариантам перевода «Барселонских обычаев» на старокаталанский язык (перевод XIII~в. по изданию Х.~Бастардеса-и-Пареры, и перевод XV~в., зафиксированный в “Constituciones de Catalunya”), отметив их различия. Данные, полученные Р.~Мейером-Германном, подтверждают наблюдение Х.~Бастардеса-и-Пареры о том, что переводчик текста рукописи K из уважения к законодательной, почти сакральной в его понимании природе текста «Барселонских обычаев» постарался сохранить синтаксические особенности оригинала и очень чётко следовал последнему, иногда на уровне дословного перевода и калькирований\footnote{\emph{Idem. }Zur Emergenz katalanischer Syntax... S. 185.}.

В своей работе «Некоторые заметки о составлении “Барселонских обычаев”» М.~Майер-Оливэ рассмотрел латинский текст Барселонских обычаев преимущественно с филологической точки зрения и сосредоточился на его языковой форме и текстуальной структуре\footnote{\emph{Mayer Olivé~M}\emph{.} Algunas notas a propósito de la compilación de los "Usatges" de Barcelona. // Estudios de Latín Medieval Hispánico: Actas del V Congreso Hispánico de Latín Medieval, Barcelona, 7--10 de septiembre de 2009. Barcelona, 2012. P.~995--1002.}. Исследователь проанализировал лексику, синтаксис и стилистику латинского оригинала, а также варианты чтения и возможные этапы редактирования, опираясь на инструментарий классической текстологии. Правовой и институциональный контекст составления «Барселонских обычаев» при этом во многом отошёл на второй план и послужил фоном для более точной реконструкции истории текста и его компиляции. Таким образом, вклад М.~Майера Оливэ в изучение «Барселонских обычаев» заключался в том, что он уточнил латинскую текстуальную традицию, в которой они возникли.

В современной западноевропейской историографии наблюдается устойчивый и последовательно углубляющийся интерес к комплексному изучению многоязычных письменных практик Средневековья, где особое внимание уделяется механизмам взаимодействия латинского и романских языков в правовых текстах --- как в сфере повседневной документации, так и в более формализованных корпусах нормативных актов. Исследователи всё чаще рассматривают перевод и языковые трансформации как сложные социокультурные процессы (формы переработки авторитета, присвоения правовых смыслов и переговоров о статусе правовой нормы), а не только технические операции межъязыковой коммуникации или механический перенос терминологии. Методологический инструментарий современных исследований включает междисциплинарные подходы, сочетающие лингвистический, антропологический и социально-исторический анализ с дипломатикой, палеографией и источниковедческой критикой, а также обращение к микроисторическим и социолингвистическим сюжетам. При этом каталонский материал остается недостаточно изученным с точки зрения трансформации правовых текстов и языковых практик.

Однако не появилось комплексных историко-лингвистических исследований, которые бы системно и на репрезентативном источников анализировали механизмы перевода, адаптации и институциональной циркуляции правовых текстов в каталонском регионе XII--XIII~вв., включая вопросы выбора языка, переработки формуляра и переноса правовых категорий. Восполнение этого пробела требует интеграции методов лингвистического анализа с социально-исторической реконструкцией медиаторов и практик письма, а также разработки сопоставимых критериев периодизации и картирования процессов «вернакуляризации» права.

Обращение к социолингвистике позволяет выйти за пределы традиционной историко-правовой оптики и рассматривать право как форму социальной коммуникации. Многоязычие, перевод и соотношение устного и письменного приобретают статус ключевых факторов, определяющих доступ к праву и механизмы его легитимации. Показано, что правовая культура Средневековья разворачивалась в условиях постоянного языкового посредничества между латынью и различными вернакулярами. Это открывает возможность по-новому интерпретировать источники, связанные с каталонским обычным правом, обращая внимание не только на содержание норм, но и на их языковую оболочку. Таким образом, социолингвистическая перспектива становится не вспомогательной, а структурообразующей для данной диссертации.

\subsection*{§3. Отечественная традиция изучения права Пиренейского полуострова в Средние века}
\addcontentsline{toc}{subsection}{§3. Отечественная традиция изучения права Пиренейского полуострова в Средние века}

Существует несколько развернутых историографических очерков\footnote{Среди таких обзоров стоит упомянуть, например, очерк \emph{В.А. }\emph{Ведюшкина}\emph{, И.С.~Пичугиной и Г.А}. \emph{Поповой }в коллективной монографии «История Испании» 2012 года: \emph{Ведюшкин}\emph{ В.А., Пичугина И.С., Попова Г.А}. Введение // История Испании под ред. \emph{В.А.~Ведюшкина}, \emph{Г.А.~Поповой}, М.: Индрик, 2012. T.~1. C.~11--17. А также краткий очерк \emph{И.С.~Филиппова} (что важно, посвящённый изучению именно каталонской истории в отечественной традиции): \emph{Filippov}\emph{ I. }Medieval Catalonia in Russian Research // Catalunya i Europa a través de l’Edat Mitjana. 2002. P.~113--122. В другом историографическом обзоре, несмотря на его обстоятельность и актуальность, содержится ряд неточностей: \emph{Ауров}\emph{ О.В.} Изучение средневековой истории Испании в России и СССР: люди, эпохи, идеи // Новый исторический вестник. 2015. Т.~13. №~160. С.~174--201.} с характеристикой изучения в том числе и каталонской истории российскими исследователями, мы будем ссылаться не только на обзоры, но и непосредственно на работы, имеющие отношение к нашей проблематике. Подобные обзоры, как правило, сосредоточены на кастильском материале или рассматривают историю Испании в целом, тогда как каталонская проблематика представлена в них лишь отдельными сюжетами. Поэтому их исследовательская перспектива лишь частично пересекается с проблемным полем настоящей работы.

\subsubsection*{§3.1. Дореволюционные исследования по истории Каталонии в Средние века}
\addcontentsline{toc}{subsubsection}{§3.1. Дореволюционные исследования по истории Каталонии в Средние века}

Отечественная традиция изучения правовых аспектов каталонского средневековья берёт начало в XIX~в. в работах И.В.~Лучицкого, а именно в посвящённых пиренейским сюжетам статье и рецензии на книгу О.~Брютайльса\footnote{\emph{Лучицкий}\emph{ И.В.} Поземельная община в Пиренеях // Отечественные записки. 1883. №~9--10, 12 (270--271). C.~57--78, 467--506, 415--442; \emph{Он же.} Brutailis A. Русские рабы и рабство в Руссильоне в XIV и XV~вв. A. Brutails, étude sur l’esclavage en Roussillon du XIII au XVII s. 1886: Рец. Киев: Унив. тип, 1886. 30 c.}. Наибольший интерес в контексте нашего диссертационного исследования представляет вторая работа, посвящённая институту рабства и критике подхода О.~Брютайльса. Особенно важно для нас то, что уже в XIX~в. территории Каталонии (Старой и Новой) и Руссильона рассматриваются как единое пространство, что и О.~Брютайльс, и И.В.~Лучицкий обращаются к тем же самым источникам, что и мы в настоящем исследовании, --- сборникам норм обычного права, составленным в XI--XIII~вв. Однако дореволюционных исследователей история права интересовала как прикладная дисциплина.

Ещё одним исследователем дореволюционного периода стал ученик И.В.~Лучицкого В.К.~Пискорский. Этот автор стал для отечественной историографии первопроходцем в изучении истории Испании в целом и Каталонии в частности. Богатый материал для его работ дала длительная командировка в Испанию, в ходе которой ему удалось познакомиться с рукописными собраниями Каталонии. В.К.~Пискорский был хорошо знаком с текстом «Барселонских обычаев», активно использовал их в своих работах в качестве одного из основных источников, например, в монографии «Крепостное право в Каталонии в средние века»\footnote{\emph{Пискорский}\emph{ В.К.} Крепостное право в Каталонии в средние века. Киев: Тип. Ун-та св. \emph{Владимира Н.Т.}~Корчак-Новицкого, 1901.}. Для этого исследования «Барселонские обычаи» стали одним из главных источников, причём впервые они были рассмотрены с исторической точки зрения. Были также предприняты попытки перевода отдельных установлений на русский язык. Не менее известна другая работа В.К.~Пискорского «Вопрос о значении и происхождении шести "дурных обычаев" в Каталонии»\footnote{\emph{Он же.} Вопрос о значении и происхождении шести «дурных обычаев» в Каталонии: с прил. неизд. Документов. Киев: тип. Ун-та св. \emph{Владимира Н.Т.}~Корчак-Новицкого, 1899.}, где ему удалось не только привести довольно обстоятельный для своего времени и первый на русском языке обзор публикаций, посвящённых «Барселонским обычаям», но и коснуться вопроса о распространении правовых норм, зафиксированных в обычаях на территории сначала Новой Каталонии, а затем и Руссильона\footnote{Там же. С.~5--7.}. Вполне естественно, что изучение социально-экономической и правовой истории Каталонии не было единственной целью зарубежной командировки В.К.~Пискорского, о чём он писал в отчёте о ней. Практически параллельно с исследованиями каталонской истории он занимался монографией о кастильских кортесах\footnote{\emph{Он же.} Кастильские кортесы в переходную эпоху от Средних веков к Новому времени: (1188--1520): исследование. Киев: Тип. Императ. университета св. Владимира, 1897. 212 c.}, подготовкой обобщающих заметок по современной ему испанской историографии\footnote{\emph{Он же.} Отчет о заграничной командировке приват-доцента университета св. Владимира В. Пискорского за 1896 и 1897 годы // Журнал Министерства народного просвещения. 1898. №~12 (320). C.~19--47.}. В.К.~Пискорский написал также общую работу «История Испании и Португалии»\footnote{\emph{Он же.} История Испании и Португалии. 2-е изд., доп-е изд., Санкт-Петербург: Акц. о-во Брокгауз-Ефрон, 1909.}, которая в силу обзорного характера значительно беднее по сравнению с его исследованиями, посвящёнными собственно каталонским сюжетам. После безвременной кончины В.К.~Пискорского, работавшего в Киевском и Казанском университетах, у него не осталось учеников, которые бы продолжили занятия сюжетами, связанными с каталонским правом.

Однако стоит учесть, что ограниченный доступ к источникам и отсутствие устойчивых научных центров изучения средневековой истории Пиренейского полуострова существенно сдерживали формирование самостоятельной исследовательской традиции в дореволюционной России. Тем не менее этот этап в развитии отечественной испанистики важен, поскольку именно на нём наметился устойчивый интерес к каталонской истории и праву.

Таким образом, в изучении каталонской истории в целом и в особенности её правовых аспектов случился перерыв, затянувшийся до середины 1930-х гг., когда оно было продолжено исследователями так называемой ленинградской школы (И.В.~Арского, А.Е.~Кудрявцева).

\subsubsection*{§3.2. Советская историография по истории Каталонии в Средние века}
\addcontentsline{toc}{subsubsection}{§3.2. Советская историография по истории Каталонии в Средние века}

Советская историография, посвящённая проблемам испанской истории, сформировалась только к 1930-м гг. Это произошло по ряду причин, детальное рассмотрение которых не входит в задачи нашего исследования, однако следует заметить, что не удалось сохранить преемственность по отношению к дореволюционной традиции, например к трудам В.К.~Пискорского. В 1930-е гг. интерес к Испании, её истории и культуре стремительно возрос благодаря сближению республиканцев с советским правительством и последующей поддержкой Советским Союзом в ходе Гражданской войны 1936--1939~гг.

Если работы А.Е.~Кудрявцева\footnote{\emph{Кудрявцев~А.Е.} Основные проблемы изучения средневековой Испании // Культура Испании. М.: Издательство Академии наук СССР, 1940.C.~42--68; \emph{Он же}. Испания в Средние века. Л.: Соцэкгиз, 1937. 256~с.} имеют для нас прежде всего историографическую ценность, то научному наследию И.В.~Арского стоит уделить больше внимания в контексте основной части настоящего диссертационного сочинения. Его квалификационная работа «Аграрный строй Каталонии IX--XII веков», защищённая в ЛГУ в 1937~г., была посвящена социально-экономическим аспектам каталонской истории раннего и высокого Средневековья и встретила неоднозначную реакцию\footnote{\emph{Добиаш}\emph{-Рождественская О.А.} Выступление на защите кандидатской диссертации \emph{И.В.~Арского} «Аграрный строй Каталонии IX--XII~вв.» (ЛГУ, 29.XI.1937~г.). // Культура западноевропейского Средневековья: научное наследие. М.: Наука, 1987. C.~239--244.}. На материалах этой работы была основана единственная монография И.В.~Арского «Очерки по истории средневековой Каталонии до соединения с Арагоном (VIII--XII века)»\footnote{\emph{Арский И.В.} Очерки по истории средневековой Каталонии до соединения с Арагоном (VIII---XII века). Л.: Издательство ЛГУ, 1941. 133 с.}. В ней целый раздел посвящён «Барселонским обычаям»\footnote{\emph{Он же. }Указ. соч. С.~25--48.} как источнику по социально-экономической истории данного региона. Этот отечественный исследователь руководствовался представлениями о «Барселонских обычаях», существовавшими в научном сообществе 1930-х гг. и с сомнением относился к их датировке XI~в. Буквально на первых страницах монографического исследования И.В.~Арский упоминает о своих контактах с французскими и каталонскими коллегами, а именно с Ж.~Кальметтом\footnote{\textbf{Жозеф }\textbf{Кальмет}\textbf{ (1873 -- 1952~гг.)} --- французский историк средних веков, профессор университета г.~Тулузы.} и Ф.~Вальс-Табернером\footnote{\textbf{Ферран Вальс-}\textbf{Табернер}\textbf{ (1888 -- 1942)---} каталонский историк Средних веков. В этой связи нам кажется важным сделать небольшое отступление. Ф. Вальс-Табернер симпатизировал советской России в целом, в 1928--1929~гг. он даже посетил с туристическим визитом Ленинград и Москву, о чём свидетельствуют опубликованные им очерки в газете “\emph{La Vanguardia}\emph{”}. Позднее опубликованы отдельным изданием: \emph{Valls i Taberner F.} Un viatger català a la Rússia de Stalin (1928) / ed. E. Zurawka, Barcelona: PPU, 1985. 135 p. Говорить о раннем знакомстве с \emph{И.В.~Арским} в этот период не приходится, потому что лишь к 1928~г. относится зачисление последнего в ЛГУ. К сожалению, у нас не было возможности получить данные из эпистолярного архива Ф.~Вальса-Табернера.}. Судя по тексту монографии И.В.~Арского, ему удалось наладить переписку и пересылку научной литературы с Ф.~Вальс-Табернером, поскольку он цитирует фрагмент из ответа ему со стороны каталонского учёного\footnote{\emph{Арский И.В.} Указ. соч. С.~42.}. Важно понимать, что исследовательский фокус был направлен именно на реконструкцию реалий (в том числе процессуальных особенностей судебного процесса), то есть явно выраженной в источнике информации. Дальнейшие исследования И.В.~Арского были прерваны Великой Отечественной войной: он ушёл добровольцем на фронт и погиб в ноябре 1941~г.

Политическую и социально-экономическую историю Каталонии, в том числе с привлечением «Барселонских обычаев», изучала Л.Т.~Мильская\footnote{\emph{Мильская}\emph{ Л.Т.} Очерки из истории деревни в Каталонии X--XII~вв. М., 1962. 153 c; \emph{Она же.} Власть графов-королей Барселонского дома и управление доменами короны // Власть и политическая культура в средневековой Европе Ч. 1. М., 1992. C.~100--105.}. Её работы по истории Каталонии были выполнены в основном в русле социально-экономических исследований, традиционных для середины ---второй половины XX~в. Поскольку изначальной её специализацией была аграрная история Германии, одним из основных методов, которыми пользовалась Л.Т.~Мильская в исследованиях по социально-экономической истории Каталонии, был компаративный. Однако Л.Т.~Мильская опубликовала целый ряд работ в сборниках и коллективных монографиях, посвящённых также институциональной и политической истории Каталонии\footnote{К их числу относятся следующие статьи: \emph{Мильская}\emph{ Л.Т.} Аграрное развитие Каталонии и окружающая среда (XII -- начало XIII~в.) // СВ. 1981. №~44. C.~117--135; \emph{Она же.} К вопросу о структуре господствующего класса в Каталонии X---XII~вв. // СВ. 1984. №~47. C.~21--28; \emph{Она же.} Пригородное землевладение арабов и Реконкиста в Каталонии XII века // Проблемы испанской истории. / под ред. \emph{С.П.~Пожарской}, М.: Наука, 1984. C.~158--165; \emph{Она же.} К вопросу о судьбах общины в Каталонии X--XII~вв. // СВ. 1985. №~48. C.~38--46; \emph{Она же.} Феодальная собственность и государственная власть в Каталонии в эпоху завершения Реконкисты // Проблемы испанской истории / под ред. \emph{С.П.~Пожарской}. М.: Наука, 1992. C.~160--165.}. Несмотря на то что она в своих работах систематически привлекала и доступный актовый материал, и исторические сочинения, «Барселонские обычаи» оставались одним из основных источников в её научных трудах. Особый интерес представляют для нас наблюдения исследовательницы относительно лингвистических особенностей каталонского актового материала, даже того, что был доступен в публикациях\footnote{\emph{Она же.} Указ. соч. 1962. С.~8.}, однако, например, о языковых особенностях нормативных текстов Л.Т.~Мильская практически полностью умалчивает. Стоит также упомянуть, что она довольно высоко оценивала проделанное И.В.~Арским исследование с поправкой на сопутствующие условия и более современные её работе труды каталонских и испанских историков, например, Е.~Родон-Бинуэ\footnote{Там же. С.~7--8.}.

В советской испанистике заметное место занимали работы А.Р.~Корсунского, в том числе его обобщающий труд «История Испании IX--XIII веков: Социально-экономические отношения и политический строй Астуро-Леонского и Леоно-Кастильского королевства». Несмотря на преимущественно североиспанский материал исследования, положения автора о составе и формировании придворной курии позволяют точнее интерпретировать соответствующие нормы правового памятника и обстоятельства его составления\footnote{\emph{Корсунский А.Р.} История Испании IX--XIII веков: Социально-экономические отношения и политический строй Астуро-Леонского и Леоно-Кастильского королевства. М.: Высшая школа, 1976. 238 c. Среди других его работ см. \emph{Корсунский А.Р.} Готская Испания: Очерки социально-экономической и политической истории. М.: Изд-во Моск. ун-та, 1969. 326 c.}. А.Р.~Корсунский уделял значительное внимание устройству средневекового общества, и эта проблематика получила продолжение в работах его многочисленных учеников\footnote{\emph{Он же.} История Испании IX--XIII веков… С.~87.}.

Таким образом, в советской традиции право и обычай часто интерпретировались сквозь призму классовых отношений, социальной борьбы и эволюции феодальных институтов. Доступ к зарубежным источникам и научным дискуссиям был ограничен до 1980-х гг., что сказалось на полноте исследований. Тем не менее в советский период были поставлены важные вопросы о взаимодействии обычного и писаного права, роли местных общин и городов, которые сохраняют актуальность по сей день.

\subsubsection*{§3.3. Современная отечественная историография права Пиренейского полуострова}
\addcontentsline{toc}{subsubsection}{§3.3. Современная отечественная историография права Пиренейского полуострова}

Граница между советским и постсоветским периодом в отечественной историографии весьма условна и практически неразличима. Многие исследователи продолжили свои штудии, начатые в советское время, в новых условиях, получив доступ к новым источникам и возможность более тесного взаимодействия с зарубежными коллегами. Период 1980--2000-х гг. можно назвать одним из самых продуктивных для отечественной медиевистики в целом и для изучения Пиренейского Средневековья в частности. В~историографическом поле появлялись сюжеты, близкие к проблематике данной диссертации, --- правовая культура, мультилингвизм, переводы правовых текстов.

Среди учеников А.Р.~Корсунского стоит упомянуть О.И.~Варьяш, С.Д.~Червонова. Несмотря на то что оба исследователя не занимались каталонским материалом, их наблюдения, сделанные при работе с текстами законодательного характера, чрезвычайно важны для данной диссертационной работы.

Область научных интересов О.И.~Варьяш была чрезвычайно широка: от аграрной истории Астуро-Леонского королевства\footnote{Этому направлению исследований был посвящён ряд её ранних статей и диссертация на соискание учёной степени кандидата исторических наук. \emph{Варьяш}\emph{ О.И.} Колонизация и крестьянство в Леоне и Кастилии в IX--XI веках // Проблемы испанской истории: сборник. 1979. C.~322--333. \emph{Она же.} О положении сервов и либертинов в Леонском королевстве в X--XI~вв. // СВ. 1980. №~43. C.~207--222.}, португальской истории до антропологии средневекового права (которая во многом стала отправной точкой для нашей методологии). Антропология права занимала заметное место в её статьях, но отдельного обобщающего монографического исследования она подготовить не успела. Большинство её статей (в том числе и не публиковавшиеся ранее\footnote{Например, \emph{Она же.} История права и антропология права // Пиренейские тетради: право, общество, власть и человек в средние века. М.: Наука, 2006. С.~27--30; \emph{Она же.} Юридическая культура португальского двора XIV~в. (corte --- “двор”, corte --- “суд”) // Там же. С.~64--68.}), посвящённых сюжетам, связанным с антропологией права и анализом средневековых законодательных текстов, были посмертно собраны и опубликованы в монографии, озаглавленной «Пиренейские тетради»\footnote{\emph{Варьяш}\emph{ О.И.} Пиренейские тетради: право, общество, власть и человек в средние века. М.: Наука, 2006. 451~c.}. Стоит отметить, что Ольга Игоревна была основателем и научным руководителем круглых столов по истории средневекового права «Право в средневековом мире», итоги которых позднее были опубликованы в одноименных сборниках «Право в средневековом мире»; проект был продолжен её учениками и  в 2000-х гг.\footnote{Право в средневековом мире~: сборник статей / под ред. \emph{О.И.~Варьяш}. М.: Изд-во ИВИ РАН, 1996; Право в средневековом мире~: сборник статей / под ред. \emph{О.И.~Варьяш}, \emph{И.И.~Варьяш}. СПб.: Алетейя, 2001; Право в средневековом мире: сборник статей / под ред. \emph{И.И.~Варьяш}, \emph{Г.А.~Поповой}. М.: ИВИ РАН, 2007; Право в средневековом мире~: сборник статей / под ред. \emph{И.И.~Варьяш}, \emph{Г.А.~Поповой}. М.: ИВИ РАН, 2008; Право в средневековом мире~: сборник статей / под ред. \emph{И.И.~Варьяш}, \emph{Г.А.~Поповой}. М.: ИВИ РАН, 2009; Право в средневековом мире: сборник статей / под ред. \emph{Г.А.~Поповой}. М.: ИВИ РАН, 2010.}. Концептуальной по содержанию является её статья «Язык средневекового права»\footnote{Там же. С.~86--90.}, в которой содержится ряд наблюдений Ольги Игоревны, сделанных в рамках работы с текстами португальских форалов. Эти замечания позволили нам поставить существование двуязычных правовых сборников в широкий пиренейский контекст и прояснить вопросы, связанные с существенными различиями между старокаталанскими и латинскими версиями таких памятников\footnote{Там же. С.~90.}.

Ещё одним исследователем-испанистом и учеником А.Р.~Корсунского, как мы уже ранее упомянули, являлся С.Д.~Червонов. Он также специализировался на работе с источниками правового характера, однако сферой его интересов вполне в духе марксистского направления в историографии была преимущественно экономическая история без обращения к правовой антропологии, да и к собственно правовой истории в целом\footnote{Об этом можно судить по той проблематике, которая наиболее интересовала \emph{С.Д.~Червонова}: \emph{Червонов}\emph{~С.Д.} Испанский средневековый город. / под. ред. \emph{О.В.~Аурова}. Москва: б.и., 2005. С.~34--35. См. также ряд его статей в журналах и сборниках: \emph{Он же.} Торговля в испанском городе XII--XIII веков (по материалам фуэрос) // Проблемы испанской истории: сборник. 1984. C.~146--157; \emph{Он же.} Земледелие и скотоводство в кастильском городе Куэнка в XII~в. // Проблемы истории Античности и Средних веков. М., 1982. С.~56--71.}. Для нас важен прежде всего опыт данного исследователя по работе с текстами правового характера.

Учеником С.Д.~Червонова можно считать В.А.~Кучумова, подготовившего диссертацию под руководством Е.~Гутновой после гибели научного руководителя. Диссертационное исследование В.А.~Кучумова было посвящено проблеме становления сословно-представительных собраний в Арагонской Короне в XII--XV~вв. и основывалось на протоколах заседаний кортесов, принятых ими постановлениях и дополнениях к уже существовавшим на тот момент правовым памятникам\footnote{Статьи, опубликованные по теме диссертации «Становление сословно-представительной монархии в Арагоне и Каталонии в XII--XV~вв.»: \emph{Кучумов В.А.} Кортесы Арагона и Каталонии: генезис и специфика / под ред. \emph{С.П.~Пожарской}. М.: Наука, 1992. C.~146--159. \emph{Он же.} Постоянные депутации арагонских и каталонских кортесов в XIII--XV~вв. Из истории европейского парламентаризма / под ред. \emph{С.П.~Пожарской}, \emph{О.И.~Варьяш}, Институт всеобщей истории, М.: ИВИ РАН, 1996. C.~21--39.}.

Ещё одним из учеников А.Р.~Корсунского, который интересовался историей Каталонии в ранний период, является И.С.~Филиппов. В его фундаментальную монографию «Средиземноморская Франция в раннее средневековье. Проблема становления феодализма» вошли отдельные параграфы, основанные на каталонском материале\footnote{\emph{Филиппов И.С.} Средиземноморская Франция в раннее средневековье: Проблема становления феодализма. М.: Скрипторий, 2000. С.~78, 492, 497--499.}. Отдельный интерес представляли сведения, связанные с исследованием особенностей климата, ландшафта и исторической географии Каталонии. Кроме того, непосредственно история права также входит в круг научных интересов Игоря Святославовича\footnote{См. например: \emph{Он же.} Римское право и Вульгата: право, закон, справедливость // Древнее право --- Ius Antiquum. 1999. №~4. С.~144--166; \emph{Он же.} Собственность между правом и экономикой~// СВ. 2020. Т.~81. №~1. С.~85--115.}.

В числе работ, которые задали нашему исследованию историко-антропологический ракурс, монографии и статьи И.И.~Варьяш\footnote{\emph{Варьяш}\emph{ И.И.} Правовое пространство ислама в христианской Испании XIII--XV~вв. 2-е изд., М.: Эдиториал УРСС, 2012; \emph{Она же.} Сарацины под властью арагонских королей: исследование правового пространства XIV~в. СПб.: Евразия Clio, 2016; \emph{Она же.} Мусульманская Европа. Сигналы идентичности. СПб.: Наука, 2020; \emph{Она же.} Обычное право в контексте взаимодействия королевского и мусульманского права в Арагонской Короне XIV~в. // Вестник Московского университета. Серия 8: История. 2016. №~5. С.~3--17.}. Её исследования представляют собой образец использования историко-антропологического подхода посредством анализа языковых особенностей, зафиксированных в источниках, и позволяют представить Арагонскую Корону как мультиязычное, мультикультурное пространство. Работы Ирины Игоревны основаны не только на трудах по фикху, записях обычного права, королевского законодательства, но и на актовом материале, связанных с судопроизводством, что придает им антропологический ракурс. Так, например, обращалась она и к вопросу о записи и способах принесения судебных клятв в правовой практике XIV--XV~вв.\footnote{Там же. 101.}. Примечательно и внимание Ирины Игоревны к «Обычаям Тортосы» и «Фуэро Валенсии», к текстам которых она в силу проблематики своего исследования обращается, конечно, эпизодически\footnote{\emph{Она же.} Сарацины под властью арагонских королей.... С.~5, 210--211.}.

Отдельный пласт исследований представляют собой работы по истории права иноконфессиональных групп на Пиренейском полуострове, написанные в 2000--2020-х гг., выполненные ученицей О.И.~Варьяш Г.А.~Поповой\footnote{\emph{Попова Г.А.} Сунна христиан, Lex Visigothorum, Книга судей, Фуэро Хузго: новая жизнь вестготского права // Historia animatа. Ч. 1. М., 2004. С.~142--155; \emph{Она же.} Кто говорил на арабском языке в христианском Толедо? (сферы использования арабского языка в Толедо XII--XIII~вв.) // Одиссей. Человек в истории. М.: Наука, 2008. С.~165--177; \emph{Она же.} Судья судье рознь? Судебные полномочия и их обладатели в Толедо XII--XIV~вв. // Право в средневековом мире. 2010. М., 2010. С.~102--111; \emph{Она же.} Правовая норма в актах Толедо XII--XIII~вв. // СВ. 2015. Т.~76. №~3--4. С.~320--332; \emph{Она же.} Процесс составления городских сводов права в Кастилии XII--XIII~вв.: норма и практика // Электронный научно-образовательный журнал История. 2017. Т.~8. №~6 (60). Доступ для зарегистрированных пользователей. URL:~\url{https://history.jes.su/s207987840001891-9-1} (дата обращения: 13.10.2024).} по материалам обычного права и королевского законодательства, а также актам в Толедском королевстве XI--XIII~вв. Работы Галины Александровны посвящены различным аспектам осмысления и бытования средневекового права (как основанного на более ранних правовых текстах, например, «Вестготской правде», так и имеющего в основе обычай, право общины или право королевское) на территории Толедского королевства в интересующий нас период, что позволяет обращаться к ним в качестве сравнительного материала для компаративного исследования. Они органично вписываются в рамки направления, заданного О.И.~Варьяш.

Среди учеников И.И.~Варьяш следует назвать М.А.~Астахова, исследующего дипломатические отношения Арагонской Короны и Гранадского эмирата в XIV--XV~вв., в том числе клятвенные практики и организацию арагонской дипломатической службы\footnote{\emph{Астахов М.А.} Клятвы в дипломатических отношениях Арагонской Короны и Гранадского эмирата в XIV--XV~вв. // Одиссей: человек в истории. Ритуалы и религиозные практики иноверцев во взаимных представлениях / отв. ред. \emph{С.И.~Лучицкая}. 2017. С.~46--63. См. также его другие работы, например историографический обзор: \emph{Он же. }Гранадский эмират в отечественной историографии // Электронный научно-образовательный журнал «История». 2017. Т.~8. №~10~(64). Доступ для зарегистрированных пользователей. URL: \url{https://history.jes.su/s207987840001969-4-1/} (дата обращения: 13.10.2024); \emph{Он же.} Арагонская дипломатическая служба и арагонско-гранадские отношения в эпоху поздней Реконкисты // Электронный научно-образовательный журнал «История». 2016. T.~7. Вып.~8 (52). Доступ для зарегистрированных пользователей. URL:~\url{https://history.jes.su/s207987840001605-4-1/} (дата обращения: 13.10.2024).}.

К иному направлению этой исследовательской традиции относятся работы Д.П.~Сафроновой, посвящённые правовой и документальной практике Леонского королевства X--XI~вв. В центре её внимания находятся использование документов, соотношение устной речи и письменной фиксации, а также проблема подлинности грамот\footnote{Например, см. \emph{Сафронова Д.П.} Документы на ассамблеях в Леонском королевстве X--XI веков // СВ. 2018. Т.~79. №~2. С.~86--106; \emph{Она же.} Устная речь в правовой и документальной практике Леонского королевства X--XI веков // Диалог со временем. 2020. №~72. С.~124--137; \emph{Она же.} Фальшивые грамоты в леонской правовой практике X--XI веков // История: переводить, понимать, оценивать (к юбилею \emph{М.А.~Юсима}). 2021. С.~121--140.}. Эти сюжеты рассматриваются преимущественно в историко-антропологической перспективе и с особым вниманием к языку источников. При сравнительном анализе исследовательница обращается также к материалам каталонской правовой традиции\footnote{См. например, \emph{Она же.} Устная речь в правовой и документальной практике Леонского королевства X--XI веков // Диалог со временем. 2020. №~72. С.~125.}.

К микроисторическому и историко-антропологическому направлению в изучении средневекового права относятся работы О.И.~Тогоевой, опубликованные в конце 1990-х--2010-х гг.\footnote{\emph{Тогоева}\emph{ О.И.} Формы судебной клятвы во Франции XIV--XV~вв. // Право в средневековом мире. 2009. C.~155--166; \emph{Она же.} «Истинная правда»: языки средневекового правосудия. М.: Наука, 2006. 333 c.; \emph{Она же.} Дела плоти. Интимная жизнь людей Средневековья в пространстве судебной полемики. М.; СПб.: Центр гуманитарных инициатив, 2018. 350 c.}. Ольга Игоревна специализируется на истории позднесредневековой Франции, поэтому в поле зрения её исследований зачастую попадают важные для нас историко-правовые сюжеты. Этого автора отличает внимание к антропологическим аспектам и интерес к иконографическим особенностям изображения правовых практик в позднее Средневековье. Каталонская правовая традиция, в отличие, например, от кастильской или португальской имела гораздо больше параллелей и общности с южно-французской традицией. Как пишет О.И.~Тогоева в заключении к своей монографии «“Истинная правда”: языки средневекового правосудия», определила свою основную задачу следующим образом: «проследить некоторые особенности средневекового правосознания, понять, как в то время воспринимались уголовный процесс, тюремное заключение, система наказаний, институт обязательного признания»\footnote{\emph{Она же.} «Истинная правда»: языки средневекового правосудия. М.: Наука, 2006. С.~299.}.

Методологический интерес для нашего исследования в области применения количественных методов для анализа структуры и содержания текстов правового характера представляют работы Е.Н.~Кирилловой\footnote{\emph{Кириллова Е.Н.} Структура ремесленных уставов по «Книге ремёсел» Этьена Буало // СВ. М., 1997. №~60. С.~89--106; \emph{Она же.} «Пошлины» и «обычаи»: о термине «Кутюмы» в Парижской «книге ремесел» // СВ. 2013. Т.~74. №~1--2. С.~199--235; \emph{Она же. }Количественные данные в историческом исследовании: французские ремёсла в XIII и XVIII~вв. // Проблемы исторического познания. 2016. №~10. С.~267--291. \emph{Она же.} История ремесла во Франции XIII--XVIII~вв.: стать мастером. СПб.: Евразия, 2019. 592~с. и др.}, выполненные на французском материале. Основной вид источников, с которыми работает Екатерина Николаевна, --- это статуты французских ремесленных корпораций позднего Средневековья.

Стоит также отметить работы по истории канонического права отечественной исследовательницы Е.В.~Казбековой\footnote{\emph{Казбекова Е.В.} Приемы ритмизированной прозы (cursus) в сводах декретального права и использование cursus в Новеллах Иннокентия IV (1243--1254) // СВ. 2005. №~66. С.~218--264; \emph{Она же.} Вставки и интерполяции в списках папских декретальных сводов XIII~в.: Экстраваганты // Электронный научно-образовательный журнал «История». 2013. №~6~(22). Доступ для зарегистрированных пользователей. URL:~\url{https://history.jes.su/s207987840000596-4-1/} (дата обращения: 02.05.2026); \emph{Она же.} Мнемотехнические латинские изречения в западноевропейских рукописях канонического права XIII --- начала XIV~вв. в Российской национальной библиотеке (Lat. F.V.~II. №~8 и №~24) // Романские языки и культуры: от Античности до современности. Сборник материалов VIII Международной научной конференции. М.: ООО «МАКС Пресс», 2016. С.~164--176.}. Она активно работает с рукописями текстов канонического права XIII~в., находящимися в собрании РНБ (Санкт-Петербург). Среди затронутых автором вопросов для настоящего исследования особенно значимы бытование этих текстов и проблема определения источников, использовавшихся при составлении местных сводов обычного права в Каталонии. Наряду с этим рассматриваются изучение канонического права в средневековой европейской традиции, особенности рукописной традиции и другие связанные сюжеты.

Одним из современных специалистов по истории средневековой Кастилии является О.В.~Ауров. Этот исследователь во многом является продолжателем линии С.Д.~Червонова и О.И.~Варьяш, однако защитил диссертацию в Санкт-Петербурге под руководством специалиста по истории Франции В.И.~Мажуги\footnote{Об этом он сам пишет в обзоре: \emph{Ауров}\emph{ О.В.} Указ. соч. С.~198--199.}. Научные интересы О.В.~Аурова сосредоточены на истории кастильского рыцарства и истории правовых памятников Пиренейского полуострова Раннего средневековья. Можно сказать, что автор практически полностью концентрируется на социально-экономической истории, не прибегая к сравнениям с каталонскими правовыми памятниками. В настоящее время исследователь руководит большими проектами по изданию средневековых испанских документов, среди которых стоит назвать проект по переводу и комментированию «Вестготской правды»\footnote{Который окончился изданием: Вестготская правда (книга приговоров). Латинский текст. Перевод. Исследование / отв. ред. \emph{О.В.~Ауров}, \emph{А.В.~Марей}. Москва: Русский фонд содействия образованию и науке 2012. 944 с. (Исторические источники)}.

Учеником О.В.~Аурова, продолжающим заниматься и по сей день классической историей права Пиренейского полуострова на позднем кастильском материале, является А.В.~Марей. Антропологическая проблематика в его работах заявлена скорее тематически, тогда как метод исследования остаётся классическим историко-правовым. Вслед за О.В.~Ауровым А.В.~Марей уделяет большое внимание анализу классической испанской историографии. При обращении к правовым текстам (кастильские фуэро, королевские зерцала, «Семь партид Альфонсо Мудрого» и др.) и историописанию (королевские кастильские хроники)\footnote{\emph{Марей А.В.} Язык права средневековой Испании: от Законов XII таблиц до Семи партид. М.: URSS, 2008; \emph{Он же.} Божий суд и человеческая справедливость: судебный поединок в Леоне и Кастилии XI--XIII~вв. М.: Союзник, 2011; \emph{Он же.} Дружба и доверие в испанском обществе XI--XIII~вв.: неправовые регуляторы правового поля. М.: Союзник, 2011.} автором практически игнорируется актовый материал.

Правовые тексты, созданные в эпоху Средневековья на Пиренейском полуострове, привлекали и исследователей филологического профиля. Среди таких работ можно назвать А.Н.~Иванову. Её диссертация посвящена особенностям языковой картины мира в «Семи партидах» Альфонсо X Мудрого\footnote{\emph{Иванова А.Н.} Свод законов Альфонсо Х Мудрого «Siete Partidas» (1256--1265) как отражение средневековой испанской языковой картины мира: диссертация на соискание учёной степени кандидата филологических наук. М., 2014.}. Несмотря на анализ лексического поля этого правового памятника, исследование недостаточно учитывает исторический контекст и специфику источника.

Отдельно стоит сказать о работах из области социальной антропологии, которые оказали влияние на исследовательскую оптику настоящей диссертации. Прежде всего это работы А.Н.~Кожановского, посвящённые изучению в том числе языковой идентичности жителей Пиренейского полуострова\footnote{\emph{Кожановский}\emph{ А.Н.} Быть испанцем...: традиция, самосознание, историческая память. М.: АСТ Восток-Запад, 2006; \emph{Он же.} Культурно-языковое многообразие населения Испании: теория и реальность: диссертация на соискание учёной степени доктора исторических наук. М., 2013.}. Исследования Александра Николаевича написаны на более позднем материале и каталонские данные занимают в них не так много места, однако, дают представление о современном состоянии изучения каталонского языка и истории его становления как идиома.

Ещё одним современным исследователем, работающим в рамках истории права как дисциплины юридической, является Д.Ю.~Полдников. Его первая монография «Договорные теории глоссаторов»\footnote{Нам кажется важным отметить, что по первому образованию \emph{Д.Ю.~Полдников} всё-таки является историком и выпускником-дипломником \emph{И.С.~Филиппова} (2002). Затем Полдников продолжил своё образование на юридическом факультете МГУ им. \emph{М.В.~Ломоносова} и диссертацию на соискание кандидата юридических наук, положенную в основу данной монографии, защитил уже там: \emph{Полдников Д.Ю.} Договорные теории глоссаторов. М.: Academia, 2008.} посвящена взгляду средневековых юристов-глоссаторов на договор как явление гражданского права. Легисты-глоссаторы преподавали и изучали римское и каноническое право в Болонском университете с XI по XIII~вв. и находились в авангарде западноевропейской правовой науки своего времени. Их труды представляли собой прежде всего пояснения к спорным фрагментам текстов римского права --- так называемые глоссы. Отдельные положения трактатов этих юристов-теоретиков довольно быстро распространились на территории Прованса, Окситании, Каталонии. В работе Д.Ю.~Полдникова также представлены определённые свидетельства их влияния и на общепиренейскую правовую традицию эпохи Высокого Средневековья. Д.Ю.~Полдников в своей работе обращается, прежде всего к юридическому содержанию учения «глоссаторов» о договоре в силу своей научной специализации. Этому же исследователю принадлежит первая публикация русскоязычного перевода важного в контексте составления текста «Барселонских обычаев» сочинения, а именно «Составленных Петром извлечений из римских законов»\footnote{Составленные Петром извлечения из римских законов / пер. с лат. \emph{Д.Ю.~Полдникова}, под ред. \emph{Л.Л.~Кофанова} // Древнее право. Ius antiquum. 2003. №~1 (11). С.~216--262; №~2 (12). С.~212--231; 2004. №~1 (13). С.~188--207.}. На данный момент Д.Ю.~Полдниковым написано несколько работ по теории истории права. Для нас особый концептуальный интерес представляет ряд его статей и глав в коллективных монографиях, посвящённых трактовке понятия «правовая культура» в современной юриспруденции и истории государства и права\footnote{Например, \emph{Полдников Д.Ю. }Правовая культура и правовая традиция как предмет сравнительной истории права // Историко-правовой ежегодник--2023. 2024. №~2. С.~100--123; \emph{Он же.} Особенности юридической техники западноевропейской правовой науки XIII--XV~вв. // История государства и права. 2015. №~6. С.~19--25 и др.}.

Довольно часто историки упускали из виду гражданское право и договор в целом. Одним из самых современных исследований, также посвящённых договору (но с несколько иной точки зрения), является монография П.Ю.~Шамаро «Частноправовой акт в Южной Франции X--XII~вв.»\footnote{Монография базируется на диссертации, защищённой автором в СПбИИ РАН под руководством \emph{В.И.~Мажуги} в 2023 году. В 2002 году автор окончил исторический факультет МГУ им. \emph{М.В.~Ломоносова} и защищал дипломную работу под руководством \emph{И.С.~Филиппова}. Автором был опубликован ряд статей, близких по проблематике указанной монографии: \emph{Шамаро}\emph{~П.Ю.} Частноправовой акт в Южной Франции X--XII~вв. М.: Наука, 2024. 216 с.; \emph{Он же.} Предметы-символы и письменный документ в удостоверении правовых актов в XI--XII~вв. Символы универсальные и условные // СВ. 2014. №~75 (3--4). С.~193--213; \emph{Он же.} Грамоты монастыря св. Максенция в архиве СПбИИ РАН // СВ. 2011. №~72 (3--4). С.~256--279 и др. Стоит сказать, что, несмотря на специфический характер, книга получила положительные рецензии в историческом сообществе, см. например: \emph{Филиппов И.С. }Рец.: \emph{Шамаро П.Ю.}~Частноправовой акт в Южной Франции Х--XII~вв. М.: Наука, 2024 // СВ. 2024. Т.~85. №~2. С.~226--231.}. Монография основана на обширной выборке актового материала, либо опубликованного, либо находящегося в составе архива СПбИИ РАН. Она представляет собой не только исследование дипломатических особенностей этих частноправовых актов, но и анализ средств правовой коммуникации, которые использовались при составлении таких документов. Отдельные выводы автора требуют уточнений, однако, некоторые из его наблюдений оказались весьма полезными в настоящем диссертационном исследовании.

Стоит также упомянуть работы, опубликованные в последние годы: А.А.~Белимовой\footnote{\emph{Белимова А.А.} Рождение истории в Каталонии: судьба одной сгоревшей рукописи // Диалог со временем. 2019. №~67. С.~246--262.}, посвящённые каталонскому историописанию X--XIII~вв.; И.В.~Билецкой\footnote{\emph{Билецкая И.В.} Законодательная регламентация долговых отношений иудеев и христиан в Кастилии середины XIII~в. // Клио. 2022. №~2~(182). С.~52--57; \emph{Она же.} Законодательная регламентация кредитно-долговых отношений иудеев и христиан в Кастилии XV~в. // Клио. 2023. №~2~(194). С.~48--53; \emph{Она же.} Кредитные споры в Авиле 1477~г. в контексте долговых отношений иудеев и христиан в Кастилии XV века // Человеческий капитал. 2023. №~5~(173). С.~38--48.} об экономико-правовых взаимоотношениях иудеев и христиан, хотя и выполненных на кастильском материале XIII--XIV~вв.; А.С.~Ковалевым\footnote{\emph{Ковалев А.С. }Dominium в правовой культуре пост-римской Италии: свидетельства Кассиодора // СВ. 2022. Т.~83. №~2. С.~66--88; \emph{Он же.} Церковь и ее епископы в остготской Италии: правоимущественный аспект // Вестник Православного Свято-Тихоновского гуманитарного университета. Серия 2: История. История Русской Православной Церкви. 2024. №~118. С.~9--26.} о правовых категориях в раннесредневековой Италии; Х.~Сотой\footnote{\emph{Сота Х.} Вера «простецов»: нравственно-религиозный облик крестьян в «Кантигах о Святой Марии» Альфонсо Х Мудрого // СВ. 2020. Т.~81. №~4. C.~102--128; \emph{Он же.} Король и его подданные: правовые представления в Кантигах о святой Марии Альфонсо Х Мудрого // Электронный научно-образовательный журнал «История». 2025. Т.~16. Вып.~3~(149). Доступ для зарегистрированных пользователей. URL:~\url{https://history.jes.su/s207987840035056-0-1/} (дата обращения: 02.05.2026).} о правовых представлениях в «Кантигах о святой Марии» Альфонсо Х Мудрого.

Таким образом, несмотря на многочисленные обращения отечественных исследователей к правовой истории XI--XIII~вв., в том числе и на каталонском материале, систематического исследования, предполагающего анализ правовой культуры региона и использование различных по видовой принадлежности источников, на данный момент не существует, чем во многом и объясняется актуальность представляемого нами диссертационного сочинения.

К началу XXI~в. сложилась комплексная методологическая база для изучения правовой культуры средневековой Каталонии, объединяющая достижения французской школы социальной истории, англо-американской медиевистики и классических историко-правовых исследований. Ключевым достижением западноевропейской историографии стало преодоление узко-текстологического подхода к правовым источникам и утверждение понимания широких познавательных возможностей обычного права в области исторической антропологии. Параллельно сформировался социолингвистический подход к анализу текстов, составленных сразу на нескольких языках или подвергшихся переводу. Все эти достижения позволяют рассматривать выбор языка в правовых текстах как социально значимое явление, отражающее режимы легитимации и распределение авторитета между латынью и вернакулярами.
